\documentclass[8pt,a4paper]{article}

\usepackage[UTF8]{ctex}
\usepackage[landscape,margin=1.8cm]{geometry}
\usepackage{fontspec}

\usepackage{listings}
% 让 markdown 中的 $...$ / $$...$$ 数学公式按正常 LaTeX 方式渲染
\usepackage[hybrid]{markdown}
\usepackage[dvipsnames]{xcolor}
\usepackage{fancyhdr}
\usepackage{graphicx}
\usepackage{multicol}
\usepackage{amsmath}
\usepackage{amssymb}

% 设置等宽字体 (用于代码块)
\setmonofont{JetBrainsMonoNerdFontMono-Regular.ttf}[Path=./src/]
% \setmainfont{JetBrainsMonoNerdFontMono-Regular.ttf}[Path=./src/]
\graphicspath{{./src/}{./assets/formulas/}}

% --- 为 markdown 里的 ```cpp 代码块定义别名 ---
\lstdefinelanguage{cpp}{
    language = C++
}

% --- listings 样式配置 ---
\lstset{
    language            =   C++,                                    % 语言设置
    numbers             =   left,                                   % 行号位置
    numbersep           =   6pt,                                    % 行号距离
    numberstyle         =   \tiny\color{gray},                      % 行号设置
    lineskip            =   2pt,                                    % 行间距离
    basicstyle          =   \ttfamily\scriptsize,                   % 基础文本风格
    keywordstyle        =   \ttfamily\bfseries\color{blue},         % 关键字风格
    commentstyle        =   \ttfamily\bfseries\color{ForestGreen},  % 注释风格
    stringstyle         =   \ttfamily\color{red},                   % 字符串风格
    showstringspaces    =   false,                                  % 不显示字符串中空格
    frame               =   shadowbox,                              % 边框样式
    framerule           =   0.3pt,                                  % 边框粗细
    framesep            =   2pt,                                    % 边框与代码间距
    rulesepcolor        =   \color{red!20!green!20!blue!20},        % 边框阴影颜色
    breaklines          =   true,                                   % 自动断行
}

% --- 为 shell script 单独设置 ---
\lstdefinestyle{sh}{
    language            = sh,
    keywordstyle        =   \ttfamily\bfseries\color{violet},
}

% --- 创建一个新的、方便调用的命令 ---
\newcommand{\code}[1]{
    \lstinputlisting{#1}
}
\newcommand{\shcode}[1]{
    \lstinputlisting[style=sh]{#1}
}
\newcommand{\txt}[1]{\par\noindent\footnotesize #1}
\newcommand{\inlineformula}[1]{\ensuremath{\input{#1}}}
\newcommand{\displayformula}[1]{%
    \vspace{0.2em}
    \[
        \input{#1}
    \]
}

% --- 文档信息 (用于封面) ---
\title{XCPC算法模板}
\author{至圣 \&\& Nrtusea \&\& Hiraethsoul}
\date{} % 备份版不显示日期

% --- 页眉页脚配置（备份版在页眉显示作者信息） ---
\makeatletter
\pagestyle{fancy}                       % 使用 fancy 样式
\fancyhf{}                              % 清空所有默认内容
\fancyhead[C]{至圣 \&\& Nrtusea \&\& Hiraethsoul} % 页眉中间: 作者们
\fancyfoot[C]{\thepage}                 % 页脚中间: 页码
\renewcommand{\headrulewidth}{0.4pt}    % 页眉下方横线
\renewcommand{\footrulewidth}{0pt}      % 去掉页脚横线
\makeatother

\begin{document}

% --- 封面页 ---
\makeatletter
\begin{titlepage}
    \centering
    \thispagestyle{empty}
    \vspace*{0.8cm}

    % 标题
    {\Huge\bfseries \@title\par}
    \vspace{0.75cm}
    
    \hrule

    % 副标题
    \vspace{0.35cm}
    {\Large \textit{}\par}
    \vspace{0.75cm}

    % 作者（只显示多作者信息，去掉额外联系方式与学校）
    {\Large \@author\par}
    
    \vfill

\end{titlepage}
\makeatother

% --- 目录页 ---
\newpage
\tableofcontents
\newpage

% --- 正文内容从此页开始 ---
\setlength{\parindent}{0pt}
\setlength{\columnsep}{22pt}
\begin{multicols}{2}
\pagenumbering{arabic}

\section{Prepare | 准备工作}

    \subsection{Testing | 测试}

        \subsubsection{Floyd | 传递闭包}
            \txt{\begin{itemize}
                \item Codeforces: 1421 ms
                \item Atcoder: 476 ms
                \item QOJ: 328 ms
            \end{itemize}}
            \code{src/misc/test_floyd.cpp}

        \subsubsection{Floyd with Bitset | 传递闭包 bitset 优化}
            \txt{\begin{itemize}
                \item Codeforces: 2077 ms
                \item Atcoder: 1065 ms
                \item QOJ: 1146 ms
            \end{itemize}}
            \code{src/misc/test_floyd_bitset.cpp}

        \subsubsection{Basic operations | 基本运算}
            \txt{\begin{itemize}
                \item Codeforces: 3015 ms
                \item Atcoder: 1512 ms
                \item QOJ: 1132 ms
            \end{itemize}}
            \code{src/misc/test_normal.cpp}

        \subsubsection{Float operations | 浮点运算}
            \txt{\begin{itemize}
                \item Codeforces: 1327 ms
                \item Atcoder: 1324 ms
                \item QOJ: 1233 ms
            \end{itemize}}
            \code{src/misc/test_float.cpp}

        \subsubsection{STL performance | 1e6 STL性能}
            \txt{\begin{itemize}
                \item Codeforces: 2484 ms
                \item Atcoder: 1526 ms
                \item QOJ: 1221 ms
            \end{itemize}}
            \code{src/misc/test_stl.cpp}

    \subsection{Shell script | 编译脚本}

        \subsubsection{main.sh | 编译运行}
            \shcode{src/misc/main.sh}

        \subsubsection{checker.sh | 对拍脚本}
            \shcode{src/misc/checker.sh}

        \subsubsection{checker\_spj.sh | SPJ 对拍脚本}
            \shcode{src/misc/checker\_spj.sh}

\section{Data structure | 数据结构}

    \subsection{Disjoint union set | 并查集}
        \code{src/data_structure/disjoint_union_set.cpp}

    \subsection{Disjoint union set rollback | 可撤销并查集}
        \code{src/data_structure/disjoint_union_set_rollback.cpp}

    \subsection{Disjoint union set range | 倍增并查集}
        \code{src/data_structure/disjoint_union_set_range.cpp}

    \subsection{Sparse table | ST 表}
        \code{src/data_structure/sparse_table.cpp}

    \subsection{Fenwick tree | 树状数组}
        \code{src/data_structure/fenwick_tree.cpp}

    \subsection{Segment tree | 懒标记线段树}
        \code{src/data_structure/segment_tree.cpp}

    \subsection{Segment tree persistent | 可持久化线段树}
        \txt{split(p, x, y, l, r, L, R):}
        \txt{\begin{itemize}
            \item p: 当前需要划分的根
            \item x: 划分后 [L, R) 的部分 
            \item y: 划分后剩余的部分
        \end{itemize}}
        \txt{merge(x, y, l, r):}
        \txt{\begin{itemize}
            \item x: 需要合并的树根
            \item y: 需要合并的树根
        \end{itemize}}

        \code{src/data_structure/segment_tree_persistent.cpp}

    \subsection{Segment tree beats | 区间 min/max 线段树}
        \code{src/data_structure/segment_tree_beats.cpp}

    \subsection{Segment tree line | 李超树}
        \code{src/data_structure/segment_tree_line.cpp}

    \subsection{Query tree | 线段树分治}
        \code{src/data_structure/query_tree.cpp}

    \subsection{Chtholly tree | 珂朵莉树}
        \code{src/data_structure/chtholly_tree.cpp}

    \subsection{Bit trie | 01 字典树}
        \code{src/data_structure/bit_trie.cpp}

    \subsection{Xor basis | 线性基}
        \code{src/data_structure/hamel.cpp}

    \subsection{Prefix xor basis | 前缀线性基}
        \code{src/data_structure/prefixhamel.cpp}

    \subsection{PBDS tree | PBDS 平衡树}
        \txt{inset(x): 插入一个元素 x}
        \txt{erase(it): 删除一个迭代器 it}
        \txt{order\_of\_key(x): 返回严格小于 x 的数量}
        \txt{find\_by\_order(x): 返回 cmp\_fn 排名为 x 的迭代器}
        \txt{join(x, y): 合并两个集合 x 和 y，合并到 x 上，y 被清空}
        \txt{split(x, k): 以 cmp\_fn 比较，小于等于 x 的属于当前树，其余的属于 b 树}

        \code{src/data_structure/pbds_tree.cpp}

\section{Graph | 图论}

    \subsection{Graph DP \& Tree Techniques | 图论选摘}
        \markdownInput{图论.md}


    \subsection{LCA | 最近公共祖先}
        \code{src/graph/lca.cpp}
    
    \subsection{Cycle count | 普通环计数}
        \txt{时间复杂度 $O(2^n \cdot m)$，注意需要是 0-based。}
        \code{src/graph/cycle_count.cpp}
    
    \subsection{Cut edge | 割边}
        \code{src/graph/tarjan_cut_edge.cpp}
    
    \subsection{Cut node | 割点}
        \code{src/graph/tarjan_cut_node.cpp}
    
    \subsection{SCC | 强连通分量}
        \code{src/graph/tarjan_scc.cpp}
    
    \subsection{EBCC | 边双连通分量}
        \code{src/graph/tarjan_ebcc.cpp}
    
    \subsection{VBCC | 点双连通分量 (点集)}
        \code{src/graph/tarjan_vbcc_node.cpp}
    
    \subsection{VBCC | 点双连通分量 (边集)}
        \code{src/graph/tarjan_vbcc_edge.cpp}
      
    \subsection{2-SAT | 2-SAT}
        \code{src/graph/2-sat.cpp}
    
    \subsection{Virtual tree | 虚树}
        \code{src/graph/virtual_tree.cpp}

\section{Dynamic Programming | 动态规划}

    % Dynamic Programming section generated from 动态规划.md
% 所有代码块使用 lstlisting，公式按 LaTeX 语法微调

\subsection{Subsequence DP | 子序列问题}

    \subsubsection{LIS / LDS / 不升 / 不降子序列}
        \txt{最长上升 / 下降 / 不升 / 不降子序列可用二分 + 贪心在 $O(n\log n)$ 内求出。典型做法维护一个数组 $f$，$f[k]$ 为长度为 $k$ 的子序列的最小末尾值。}

    \subsubsection{LCS | 最长公共子序列}
        \txt{给定字符串 $A,B$，常见做法有：}
        \txt{1. 经典二维 DP：$dp[i][j]$ 表示 $A$ 前 $i$ 位与 $B$ 前 $j$ 位的 LCS 长度，复杂度 $O(nm)$。}
        \txt{2. 若某一串中字符不重复，可用 LIS 化简：将另一串按该串重标并删去无关字符，在新序列上做 LIS。}
        \txt{3. bitset + DP 的 bit-parallel 做法，可降到 $O(|A|\cdot |B|/w)$。下面给出字母集为小写英文字母的 bit-parallel LCS 模板。}

        \begin{lstlisting}
using u64 = unsigned long long;

// 64 位块数表示的 bitset LCS（Hunt–Szymanski 风格 bit-parallel）
// 时间 O(|a| * ((|b| + 63) / 64)), 空间 O(26 * blocks)
int LCS(const std::string &a, const std::string &b) {
    int n = (int)a.size(), m = (int)b.size();
    if (n == 0 || m == 0) return 0;

    int W = (m + 63) / 64;                   // 每个 bitset 需要的 u64 个数
    std::vector<std::vector<u64>> pos(26, std::vector<u64>(W, 0ULL));

    // 构造 pos：pos[c] 在 b 中字符 c 出现位置的位图
    for (int j = 0; j < m; ++j) {
        int c = b[j] - 'a';
        pos[c][j >> 6] |= (1ULL << (j & 63));
    }

    std::vector<u64> dp(W, 0ULL), U(W), V(W), SUB(W);
    // 最后一块只保留到 m 的有效位
    u64 last_mask = (m % 64 == 0) ? ~0ULL : ((1ULL << (m % 64)) - 1ULL);

    for (char ch : a) {
        int idx = ch - 'a';

        // U = dp | pos[idx]
        for (int i = 0; i < W; ++i) {
            U[i] = dp[i] | pos[idx][i];
        }

        // V = (dp << 1) | 1
        u64 carry = 0;
        for (int i = 0; i < W; ++i) {
            u64 next_carry = dp[i] >> 63;
            V[i] = (dp[i] << 1) | carry;
            carry = next_carry;
        }
        V[0] |= 1ULL; // 位置 0 视作空前缀

        // SUB = U - V（多块减法，考虑借位）
        u64 borrow = 0ULL;
        for (int i = 0; i < W; ++i) {
            unsigned __int128 v_with_borrow =
                (unsigned __int128)V[i] + (unsigned __int128)borrow;
            u64 v64 = (u64)v_with_borrow;
            u64 res = (u64)((unsigned __int128)U[i] - v_with_borrow);
            SUB[i] = res;
            borrow = (v_with_borrow > (unsigned __int128)U[i]) ? 1ULL : 0ULL;
        }

        // dp = U & ~SUB
        for (int i = 0; i < W; ++i) {
            dp[i] = U[i] & ~SUB[i];
        }
        // 清除最后块中超出 m 的位
        dp[W - 1] &= last_mask;
    }

    // 统计 dp 中 1 的个数就是 LCS 长度
    int ans = 0;
    for (int i = 0; i < W; ++i) {
        ans += __builtin_popcountll(dp[i]);
    }
    return ans;
}
        \end{lstlisting}

\subsection{Partition of integers | 自然数划分问题}

    \txt{考虑 $n$ 个自然数之和为 $m$ 的方案数。设 $f(i,j)$ 表示选了 $i$ 个数总和为 $j$ 的方案数。}

    \subsubsection{可重复}
        \txt{讨论是否包含数字 $1$：}
        \txt{1. 至少包含一个 $1$，则去掉一个 $1$ 后为 $f(i-1,j-1)$。}
        \txt{2. 不包含 $1$，则让当前 $i$ 个数全部减去 $1$，和变为 $j-i$，即 $f(i,j-i)$。}
        \txt{转移为 $f(i,j) = f(i-1,j-1) + f(i,j-i)$，复杂度 $O(nm)$。}

        \txt{更高阶做法是利用五边形数定理，复杂度 $O(m\sqrt{m})$。下面是一个以五边形数为核心的划分数模板（需配合取模函数 \texttt{norm}）。}

        \begin{lstlisting}
std::vector<int> calcPartition(int n) {
    std::vector<long long> a(n), p(n);
    p[0] = 1;
    if (n > 1) p[1] = 1;
    if (n > 2) p[2] = 2;
    for (int i = 1; 2 * i + 1 < n; ++i) {
        // 五边形数 1, 5, 12, 22, ...
        a[2 * i]     = 1LL * i * (3 * i - 1) / 2;
        a[2 * i + 1] = 1LL * i * (3 * i + 1) / 2;
    }
    for (int i = 3; i < n; ++i) {
        for (int j = 2; a[j] <= i; ++j) {
            // p[n] = p[n-1] + p[n-2] - p[n-5] - p[n-7] + ...
            p[i] += ((j & 2) ? 1 : -1) * p[i - a[j]];
            norm(p[i]);  // 视需要取模
        }
    }
    std::vector<int> res(n);
    for (int i = 0; i < n; ++i) res[i] = (int)p[i];
    return res;
}
        \end{lstlisting}

    \subsubsection{不可重复}
        \txt{同样按是否包含 $1$ 分类：}
        \txt{1. 恰好包含一个 $1$，则把每个数加 $1$，相当于 $f(i-1,j-i)$。}
        \txt{2. 不包含 $1$，则同样每个数加 $1$，对应 $f(i,j-i)$。}
        \txt{于是转移为 $f(i,j) = f(i-1,j-i) + f(i,j-i)$。}

\subsection{Knapsack DP | 背包 DP}

    \subsubsection{Multiple knapsack with monotone queue | 多重背包（单调队列优化）}
        \txt{经典多重背包可以拆成若干 01 背包，或用单调队列对每个模 $w$ 的同余类优化。下面是按体积模 $w$ 分组、对每个组使用单调队列的模板。}

        \begin{lstlisting}
using i64 = long long;

auto multiBag(std::vector<std::array<int, 3>> &goods, int S) {
    // S 为总背包容量，goods: {value, weight, count}
    std::vector<i64> f(S + 1, 0);
    std::vector<int> q(S + 2, 0);

    for (auto [v, w, m] : goods) {
        auto calc = [&](int i, int j) {
            // i, j 为同余类中的两个体积，i >= j
            return f[j] + 1LL * (i - j) / w * v;
        };
        for (int up = S; up + w > S; --up) {
            int l = 1, r = 0, k = up;
            for (int x = up; x > 0; x -= w) {
                for (; k >= std::max(0LL, 1LL * x - 1LL * m * w); k -= w) {
                    while (l <= r && calc(x, k) > calc(x, q[r])) --r;
                    q[++r] = k;
                }
                f[x] = calc(x, q[l]);
                if (q[l] == x) ++l;
            }
        }
    }
    return f;
}
        \end{lstlisting}

    \subsubsection{Tree knapsack | 树上背包}
        \txt{树上背包常见于“在树上选点 / 边满足某种限制”的计数或最值题。一般在 DFS 时，用临时数组 \texttt{temp} 合并子树贡献，避免子树间的重复转移。}

        \begin{lstlisting}
template <class T>
class PackageOnTree {
    const T INF = std::numeric_limits<T>::max() / 2;

public:
    std::vector<std::vector<T>> dp;

    PackageOnTree(int n, int m, int root)
        : n(n), g(n + 1), siz(n + 1),
          infCapacity(m), root(root) {
        dp.assign(n + 1, std::vector<T>(m + 1, 0));
    }

    void addEdge(int u, int v, T w) {
        g[u].emplace_back(v, w);
        g[v].emplace_back(u, w);
    }

    void work() { dfs(root, 0); }

private:
    int infCapacity;
    int n, root;
    std::vector<std::vector<std::pair<int, T>>> g;
    std::vector<int> siz;

    // u 子树内选 cnt 个点，父边权为 w 的贡献
    T calc(int u, T w, int cnt) {
        return w * (cnt * (infCapacity - cnt)
                  + (siz[u] - cnt) * (n - infCapacity - (siz[u] - cnt)));
    }

    void dfs(int u, int fa) {
        dp[u][0] = 0;
        siz[u] = 1;
        for (auto [v, w] : g[u]) {
            if (v == fa) continue;
            dfs(v, u);
            siz[u] += siz[v];
            std::vector<T> temp(std::min(siz[u], infCapacity) + 1, 0);
            for (int j = std::min(infCapacity, siz[u]); j >= 0; --j) {
                for (int k = std::max(0, siz[v] + j - siz[u]);
                     k <= std::min(siz[v], j); ++k) {
                    temp[j] = std::max(
                        temp[j],
                        dp[u][j - k] + dp[v][k] + calc(v, w, k)
                    );
                }
            }
            for (int j = 0; j <= std::min(siz[u], infCapacity); ++j) {
                dp[u][j] = temp[j];
            }
        }
    }
};
        \end{lstlisting}

\subsection{Interval DP | 区间 DP}

    \txt{区间 DP 一般形如：}
    \[
        dp[i][j] = \min_{i \le k < j} \bigl(dp[i][k] + dp[k+1][j] + \text{cost}(i,j)\bigr),
    \]
    \txt{复杂度通常为 $O(n^3)$。常用技巧包括四边形不等式 / 决策单调性优化。}

\subsection{Monotone queue optimization | 单调队列优化}

    \txt{若 DP 形如 $dp[i] = \min\_j (dp[j] + \text{cost}(j,i))$，且对每个 $i$ 可行转移 $j$ 在某个滑动窗口内，则可用单调队列维护“候选 $j$ 的代价”，将复杂度降为线性。}

\subsection{Monge decision optimization | 决策单调性优化}

    \subsubsection{2D form}
        \txt{典型二维形式：}
        \[
            dp[i][k] = \min_{1 \le j < i} \{ dp[j][k-1] + w(j+1,i) \},
        \]
        \txt{若满足决策单调性，即最优决策点在层间单调，则可用分治或二分队列优化。}

        \paragraph{Divide and conquer优化}
        \txt{在第 $k$ 层中，用分治求出区间 $[l,r]$ 的最优决策点在 $[L,R]$ 区间内，复杂度 $O(nk\log n)$。}

        \begin{lstlisting}
template <class T>
class OptimizeDP2D {
    const T INF = std::numeric_limits<T>::max() / 2;

public:
    std::vector<std::vector<T>> dp;
    int n, K;

    OptimizeDP2D(int n, int k) : n(n), K(k) {
        dp.assign(n + 1, std::vector<T>(k + 1, 0));
        for (int j = 1; j <= k; ++j) {
            work(1, n, 1, n, j);
        }
    }

private:
    // 代价函数 w(l, r) 需按题意实现
    T calcW(int l, int r) {
        // TODO: 按题目实现
        return 0;
    }

    // 在第 dep 层，求解区间 [l,r]，决策点限制在 [L,R]
    void work(int l, int r, int L, int R, int dep) {
        if (l > r) return;
        int mid = (l + r) >> 1;
        T best = -INF;
        int opt = -1;
        for (int i = std::min(R, mid); i >= L; --i) {
            T val = dp[i - 1][dep - 1] + calcW(i, mid);
            if (val > best) {
                best = val;
                opt = i;
            }
        }
        dp[mid][dep] = best;
        work(l, mid - 1, L, opt, dep);
        work(mid + 1, r, opt, R, dep);
    }
};
        \end{lstlisting}

        \paragraph{Binary-search queue优化}
        \txt{也可用“决策区间 + 队列”维护三元组 $[j,l,r]$，在队列上二分找到每个位置的最优决策点，复杂度同样为 $O(nk\log n)$。代码与 1D 形式极为相似，此处不再赘述。}

    \subsubsection{1D form}
        \txt{一维形式：}
        \[
            dp[i] = \min_{1 \le j < i} \{ dp[j] + w(j+1,i) \},
        \]
        \txt{决策单调性出现在同一层内部，无法分治，只能使用二分队列方法维护一个凸壳队列。}

        \begin{lstlisting}
template <class T>
class OptimizeDP1D {
    const T INF = std::numeric_limits<T>::max() / 2;

public:
    std::vector<T> dp;

    explicit OptimizeDP1D(int n) : n(n) {
        dp.assign(n + 1, INF);
        dp[0] = 0;
        work();
    }

private:
    int n;

    T getW(int l, int r) {
        // TODO: 按题意实现 w(l, r)
        return 0;
    }

    T getDP(int l, int r) {
        return dp[l] + getW(l + 1, r);
    }

    void work() {
        std::vector<std::array<int, 3>> q(n + 1);
        // q[t] = {j, L, R} 表示在 [L,R] 内的最优转移点为 j
        int hh = 1, tt = 0;
        q[++tt] = {0, 1, n};
        for (int i = 1; i <= n; ++i) {
            while (hh < tt && q[hh][2] < i) ++hh;
            q[hh][1] = i;
            int k = q[hh][0];
            dp[i] = getDP(k, i);

            while (hh < tt && getDP(q[tt][0], q[tt][1]) >= getDP(i, q[tt][1])) {
                --tt;
            }
            q[tt][2] = n;
            int L = q[tt][1], R = q[tt][2];
            int pos = n + 1;
            while (L <= R) {
                int mid = (L + R) >> 1;
                if (getDP(q[tt][0], mid) >= getDP(i, mid)) {
                    R = mid - 1;
                    pos = mid;
                } else {
                    L = mid + 1;
                }
            }
            if (pos != n + 1) {
                q[tt][2] = pos - 1;
                q[++tt] = {i, pos, n};
            }
        }
    }
};
        \end{lstlisting}

\subsection{Slope trick / Convex hull optimization | 斜率优化}

    \txt{斜率优化适用于形如}
    \[
        dp[i] = \min_j \bigl(dp[j] + w_1(i) + w_2(j) + f(i)g(j)\bigr)
    \]
    \txt{的 DP。令 $y = dp[j] + w_2(j)$，$x = g(j)$，$k = -f(i)$，$b = dp[i] - w_1(i)$，可化为直线形式 $y = kx + b$，问题转化为在一堆点 $(x,y)$ 中找使得 $b$ 最优的点。若 $x$ 与 $k$ 都具备单调性，可用单调队列维护凸壳；若只有 $x$ 单调，则使用二分队列。}

    \subsubsection{Monotone queue version}

        \begin{lstlisting}
template <class T>
class ConvexHullTrick {
    const T INF = std::numeric_limits<T>::max() / 2;

public:
    std::vector<T> dp;

    explicit ConvexHullTrick(int n) : n(n) {
        dp.assign(n + 1, INF);
        dp[0] = 0;
        work();
    }

private:
    int n;

    T calcX(int x) {
        // TODO: 按题意实现
        return 0;
    }

    T calcY(int x) {
        // TODO: 按题意实现
        return 0;
    }

    T calcK(int x) {
        // TODO: 按题意实现斜率 k(i)
        return 0;
    }

    long double calcSlope(int l, int r) {
        if (calcX(r) == calcX(l)) return INF;
        return (long double)(calcY(r) - calcY(l)) /
               (long double)(calcX(r) - calcX(l));
    }

    T calcVal(int i, int j) {
        // TODO: 返回用 j 转移到 i 的代价
        return 0;
    }

    void work() {
        int hh = 0, tt = -1;
        std::vector<int> q(n + 1, 0);
        q[++tt] = 0;
        for (int i = 1; i <= n; ++i) {
            T k = calcK(i);
            while (hh < tt && calcSlope(q[hh], q[hh + 1]) <= k) ++hh;
            int j = q[hh];
            dp[i] = dp[j] + calcVal(i, j);
            while (hh < tt &&
                   calcSlope(q[tt - 1], q[tt]) >= calcSlope(q[tt], i)) {
                --tt;
            }
            q[++tt] = i;
        }
    }
};
        \end{lstlisting}

    \subsubsection{Binary-search queue version}

        \begin{lstlisting}
template <class T>
class ConvexHullTrickBS {
    const T INF = std::numeric_limits<T>::max() / 2;

public:
    std::vector<T> dp;

    explicit ConvexHullTrickBS(int n) : n(n) {
        dp.assign(n + 1, INF);
        dp[0] = 0;
        work();
    }

private:
    int n;

    T calcX(int x) { return 0; }  // TODO
    T calcY(int x) { return 0; }  // TODO
    T calcK(int x) { return 0; }  // TODO

    long double calcSlope(int l, int r) {
        if (calcX(r) == calcX(l)) return INF;
        return (long double)(calcY(r) - calcY(l)) /
               (long double)(calcX(r) - calcX(l));
    }

    T calcVal(int i, int j) { return 0; }  // TODO

    void work() {
        int hh = 0, tt = -1;
        std::vector<int> q(n + 1, 0);
        q[++tt] = 0;

        auto findOptimizePoint = [&](int L, int R, T k) {
            int ans = R;
            while (L <= R) {
                int mid = (L + R) >> 1;
                if (calcSlope(q[mid], q[mid + 1]) > k) {
                    R = mid - 1;
                    ans = mid;
                } else {
                    L = mid + 1;
                }
            }
            return q[ans];
        };

        for (int i = 1; i <= n; ++i) {
            T k = calcK(i);
            int j = findOptimizePoint(0, tt - 1, k);
            dp[i] = dp[j] + calcVal(i, j);
            while (hh < tt &&
                   calcSlope(q[tt - 1], q[tt]) >= calcSlope(q[tt], i)) {
                --tt;
            }
            q[++tt] = i;
        }
    }
};
        \end{lstlisting}

\subsection{WQS binary / weight binary search | WQS 二分}

    \txt{WQS（二分斜率）用于解决“强制选恰好 $m$ 个物品”的最优方案，前提是 $f(m)$ 对 $m$ 是凸的（一次差分单调）。}
    \txt{通过给每个被选物品减去一个惩罚权重 $\lambda$，将“选 $m$ 个”的限制吸收到代价中，变成不限制个数的简单 DP，并根据最后选的个数与目标 $m$ 的关系二分 $\lambda$。}

\subsection{Dynamic DP | 动态 DP}

    \txt{引入广义矩阵乘法 $C = A \ast B$，定义}
    \[
        C_{ij} = \max_k (A_{ik} + B_{kj})
    \]
    \txt{很多从左到右的线性 DP 可写成广义矩阵乘，与线段树结合可支持区间修改与查询。例：}
    \[
        f(i) = \max(f(i-1) + a_i, b_i),
    \]
    \txt{可写作}
    \[
        \begin{bmatrix} a_i & b_i \\ -\infty & 0 \end{bmatrix}
        \ast
        \begin{bmatrix} f(i-1) \\ 0 \end{bmatrix}
        =
        \begin{bmatrix} f(i) \\ 0 \end{bmatrix}.
    \]

\subsection{Bitmask DP | 状压 DP}

    \subsubsection{SOS DP（子集和 DP）}
        \txt{要在所有 $0 \le i < 2^n$ 上求 $\sum\_{j \subseteq i} a\_j$，朴素枚举子集复杂度 $O(3^n)$。SOS DP 的优化为：}

        \begin{lstlisting}
// 枚举子集的朴素写法
for (int s = 0; s < (1 << n); ++s) {
    for (int T = s; T; T = (T - 1) & s) {
        dp[s] += a[T];
    }
}

// SOS DP 优化：dp[i] 最终为所有子集的和
for (int j = 0; j < n; ++j) {
    for (int i = 0; i < (1 << n); ++i) {
        if (i >> j & 1) {
            dp[i] += dp[i ^ (1 << j)];
        }
    }
}
        \end{lstlisting}

\subsection{Digit DP | 数位 DP}

    \txt{数位 DP 一般按位从高到低枚举，用记忆化搜索解决。常见状态参数包括：}
    \txt{1. \texttt{pos}：当前处理的位数。}
    \txt{2. \texttt{limit}：本位是否受上界限制。}
    \txt{3. \texttt{lead0}：前缀是否全是前导零。}
    \txt{4. \texttt{last}：上一位的取值（如有相邻位约束）。}
    \txt{5. \texttt{sum} / \texttt{r} / \texttt{st}：数位和、模值、状态压缩等。}

    \begin{lstlisting}
class DigitDP {
public:
    DigitDP(int maxLen, int base)
        : a(maxLen + 1), base(base) {
        // f 按题意开维度并初始化为 -1
    }

    int solve(int L, int R) {
        return calc(R) - calc(L - 1);
    }

private:
    std::vector<int> a;
    // 具体形状视题目而定，例如 f[pos][sum][r][...]
    // std::vector<std::vector<int>> f;
    int base;

    // pos: 当前位 (从高到低)，limit: 是否受上界限制
    // lead0: 是否仍在前导零，r/sum/last/st 等按题目添加
    int dfs(int pos, bool limit, bool lead0,
            int r, int last, int sum) {
        if (pos == 0) {
            // 判断当前数是否满足题目要求
            if (/* condition */) {
                return 1;  // 或返回 sum 等属性
            } else {
                return 0;
            }
        }
        // 若不受限制且已记忆，直接返回
        // if (!limit && f[pos][...] != -1) return f[pos][...];

        int up = limit ? a[pos] : base - 1;
        int res = 0;
        for (int d = 0; d <= up; ++d) {
            bool nlimit = limit && (d == up);
            bool nlead0 = lead0 && (d == 0);
            if (nlead0) {
                // 仍在前导零
                res += dfs(pos - 1, nlimit, nlead0,
                           /* r */, /* last */, /* sum */);
            } else {
                if (/* transition is valid */) {
                    res += dfs(pos - 1, nlimit, nlead0,
                               /* new r */, /* new last */, /* new sum */);
                }
            }
        }
        // if (!limit) f[pos][...] = res;
        return res;
    }

    int calc(int x) {
        if (x < 0) return 0;
        int len = 0;
        while (x) {
            a[++len] = x % base;
            x /= base;
        }
        return dfs(len, true, true, 0, -1, 0);
    }
};
        \end{lstlisting}

\subsection{Tree DP | 树形 DP}

    \subsubsection{Rerooting DP | 换根 DP}
        \txt{换根 DP 分为“点换根”和“边换根”，常见做法有两种：}
        \txt{1. 两次 DFS + 一个数组：第一次以固定根算出子树 DP，第二次换根传递。}
        \txt{2. 两次 DFS，维护 $F(u)$（子树内答案）与 $G(u)$（子树外答案），最后将两者合并。}

        \paragraph{前后缀积 Trick}
        \txt{当转移为子树乘积形式且需要取模时，用前后缀积避免除法与逆元。例子：}
        \[
            f[u] = \prod_{v \in \text{son}(u)} (f[v] + 1).
        \]

        \begin{lstlisting}
template <class T>
class PreSufOptimizeDP {
public:
    std::vector<T> dp;

    explicit PreSufOptimizeDP(std::vector<std::vector<int>> &g)
        : g(g), dp(g.size()),
          pre(g.size()), suf(g.size()) {
        dfs1(1, 0);
        dfs2(1, 0);
    }

private:
    std::vector<std::vector<int>> g;
    std::vector<std::vector<T>> pre, suf;

    void dfs1(int u, int fa) {
        dp[u] = 1;
        pre[u].push_back(1);
        for (auto v : g[u]) {
            if (v == fa) continue;
            dfs1(v, u);
            dp[u] = dp[u] * (dp[v] + 1);
            pre[u].push_back(dp[v] + 1);
            suf[u].push_back(dp[v] + 1);
        }
        suf[u].push_back(1);
        for (int i = 1; i < (int)pre[u].size(); ++i) {
            pre[u][i] = pre[u][i] * pre[u][i - 1];
        }
        for (int i = (int)suf[u].size() - 2; i >= 0; --i) {
            suf[u][i] = suf[u][i] * suf[u][i + 1];
        }
    }

    // 原式: dp[v] = dp[v] * (dp[u] / (dp[v] + 1) + 1)
    void dfs2(int u, int fa) {
        int idx = 0;
        for (auto v : g[u]) {
            if (v == fa) continue;
            dp[v] = dp[v] * (pre[u][idx] * suf[u][idx + 1] + 1);
            ++idx;
            dfs2(v, u);
        }
    }
};
        \end{lstlisting}

\subsection{DP with after-effect | 带后效性的 DP}

    \txt{当转移存在环路或无法直接拓扑化时，DP 方程本质上是一组线性方程，可用待定系数或高斯消元求解。}

    \subsubsection{Method of undetermined coefficients | 待定系数法}
        \txt{典型应用是树上随机游走的期望问题。设 $DP(x)$ 为从 $x$ 开始到达叶子的期望步数，转移大致形如}
        \[
            DP(x) = \frac{1 + \sum\_{y \in N(x)} DP(y)}{\deg(x)},
        \]
        \txt{由于包含父子互相依赖，可设 $DP(x) = k\_x \cdot DP(fa[x]) + b\_x$，把子树的 $DP$ 形式带回父节点方程中，递推得到所有 $k,b$，最终在根处解出答案。}

    \subsubsection{Gaussian elimination | 高斯消元}
        \txt{更一般的线性 DP 可直接写成线性方程组，用高斯消元在 $O(n^3)$ 时间内解出。若系数矩阵为带状，可降到 $O(nd^2)$。}

        \txt{例如在一般图上随机游走的期望次数，可对每个点写出}
        \[
            DP(i) = \sum\_{(i,j)\in E} \frac{DP(j)}{\deg(j)} + c\_i,
        \]
        \txt{将 $DP(n)$ 固定为边界条件，在其他点上联立方程，用高斯消元即可。}



\section{String | 字符串}

    % String algorithms section generated from 字符串.md
% 所有代码块使用 lstlisting，公式按 LaTeX 语法微调

\subsection{String Hash | 字符串哈希}

    \txt{注意哈希传入的字符串都是 $1$-index。}

    \subsubsection{Ordered string hash | 有序字符串哈希}
        \txt{考虑哈希函数 $f(n)=\sum_{i=1}^{n} s[i]\cdot \text{Base}^{\,n-i}$。预处理前缀哈希 $H(i)$ 与幂数组 $Hash(i)$ 后，区间 $[l,r]$ 的哈希值为
        \[
            f[l,r] = H(r) - H(l-1)\cdot \text{Base}^{\,r-l+1}.
        \]}

        \begin{lstlisting}
struct StringHash {
    int n;
    std::vector<int> Hash;
    std::vector<int> H;
    int Base, P;

    // 默认字符串加了前导字符，s[0] 不使用
    StringHash(const std::string &s, int Base, int P)
        : n(s.size()), Hash(n), H(n), Base(Base), P(P) {
        Hash.assign(n, 0);
        H.assign(n, 0);
        Hash[0] = 1;
        for (int i = 1; i <= n - 1; ++i) {
            Hash[i] = 1LL * Hash[i - 1] * Base % P;
        }
        for (int i = 1; i <= n - 1; ++i) {
            H[i] = (1LL * H[i - 1] * Base % P + s[i]) % P;
        }
    }

    int operator()(int l, int r) const {
        return (H[r] - 1LL * H[l - 1] * Hash[r - l + 1] % P + P) % P;
    }
};
        \end{lstlisting}

    \subsubsection{Unordered string hash | 无序字符串哈希}
        \txt{考虑哈希函数 $f(n)=\sum_{i=1}^{n} h(s[i])$，其中 $h(x)$ 为预先随机的权值，只要求相同字符对应的 $h(x)$ 相同即可。}

        \begin{lstlisting}
struct StringHash {
    int n;
    std::vector<int> Hash;
    std::vector<int> H;
    int Base, P;

    // a 为 1-index 的离散值数组，max 为值域上界
    StringHash(std::vector<int> &a, int maxv, int Base, int P)
        : n(a.size()), Base(Base), P(P) {
        Hash.assign(maxv + 1, 0);
        H.assign(n, 0);
        Hash[0] = 1;
        for (int i = 1; i <= maxv; ++i) {
            Hash[i] = 1LL * Hash[i - 1] * Base % P;
        }
        for (int i = 1; i <= n - 1; ++i) {
            H[i] = (H[i - 1] + Hash[a[i]]) % P;
        }
    }

    int operator()(int l, int r) const {
        return (H[r] - H[l - 1] + P) % P;
    }
};
        \end{lstlisting}

    \subsubsection{2D ordered hash | 二维有序字符串哈希}

        \begin{lstlisting}
using u64 = unsigned long long;

struct Hash2D {
    int n, m;
    std::vector<u64> hX;                 // 行方向 Base 的幂
    std::vector<u64> hY;                 // 列方向 Base 的幂
    std::vector<std::vector<u64>> A;     // 哈希前缀和
    int P, Q;                            // 两个 Base，分别为行和列

    Hash2D(const std::vector<std::vector<char>> &a,
           int n, int m, int P, int Q)
        : n(n), m(m), hX(n + 1), hY(m + 1), A(n + 1, std::vector<u64>(m + 1, 0)),
          P(P), Q(Q) {
        hX[0] = 1;
        hY[0] = 1;
        for (int i = 1; i <= n; ++i) {
            hX[i] = hX[i - 1] * P;
        }
        for (int j = 1; j <= m; ++j) {
            hY[j] = hY[j - 1] * Q;
        }
        for (int i = 1; i <= n; ++i) {
            for (int j = 1; j <= m; ++j) {
                A[i][j] = A[i][j - 1] * Q
                        + A[i - 1][j] * P
                        - A[i - 1][j - 1] * P * Q
                        + (a[i][j] - 'a' + 1);
            }
        }
    }

    u64 query(int x1, int y1, int x2, int y2) const {
        return A[x2][y2]
             - A[x2][y1 - 1] * hY[y2 - y1 + 1]
             - A[x1 - 1][y2] * hX[x2 - x1 + 1]
             + A[x1 - 1][y1 - 1] * hX[x2 - x1 + 1] * hY[y2 - y1 + 1];
    }
};
        \end{lstlisting}

\subsection{Minimal representation | 最小表示法}

    \txt{返回一个序列字典序最小的循环同构串。常用于最小循环同构、项链问题等。}

    \begin{lstlisting}
std::vector<int> minimalString(std::vector<int> &a) {
    int n = a.size();
    int i = 0, j = 1, k = 0;
    while (k < n && i < n && j < n) {
        if (a[(i + k) % n] == a[(j + k) % n]) {
            ++k;
        } else {
            if (a[(i + k) % n] > a[(j + k) % n]) {
                i += k + 1;
            } else {
                j += k + 1;
            }
            if (i == j) ++i;
            k = 0;
        }
    }
    int start = std::min(i, j);
    std::vector<int> ans(n);
    for (int t = 0; t < n; ++t) {
        ans[t] = a[(t + start) % n];
    }
    return ans;
}
    \end{lstlisting}

    \txt{Trick：若有一堆字符串，按某种顺序拼接使得整体字典序最小，只需按比较规则 $a+b < b+a$ 排序。}

\subsection{KMP}

    \txt{对于字符串 $S$，KMP 通过前缀函数（失配函数）维护 $fail(i)$：表示以 $i$ 为结尾的前缀串的最长相等前后缀的长度。}
    \txt{模式串 $t$ 在文本串 $s$ 中的出现位置，可在 $t\#s$ 上跑一遍前缀函数，凡是 $fail(i)=|t|$ 的位置均为一次匹配。}

    \begin{lstlisting}
class KMP {
public:
    KMP() = default;

    std::vector<int> getPreFunction(std::string &s) {
        return kmp(s);
    }

    // 在 s 中找 t，返回所有匹配起始下标（0-index）
    std::vector<int> getMatchPos(std::string &s, std::string &t) {
        std::vector<int> pos;
        auto link = kmp(t + "#" + s);  // 在 s 中找 t
        for (int i = (int)t.size() + 1; i < (int)link.size(); ++i) {
            if (link[i] == (int)t.size()) {
                pos.push_back(i - 2 * (int)t.size() + 1);
            }
        }
        return pos;
    }

private:
    std::vector<int> kmp(const std::string &s) {
        int n = s.size();
        std::vector<int> link(n);
        for (int i = 1, j = 0; i < n; ++i) {
            while (j && s[i] != s[j]) {
                j = link[j - 1];
            }
            if (s[i] == s[j]) ++j;
            link[i] = j;
        }
        return link;
    }
};
    \end{lstlisting}

\subsection{Z-function（扩展 KMP）}

    \txt{对于字符串 $S$，$Z(i)$ 表示后缀 $S[i,\dots]$ 与整个串 $S$ 的最长公共前缀长度。实现时维护一个当前最右的 $Z$-box 区间 $[l,r]$ 以加速转移，约定 $Z(0)=0$。}

    \begin{lstlisting}
class ZFunction {
public:
    ZFunction() = default;

    // 单串 Z 函数
    std::vector<int> getZFunction(const std::string &s) {
        return zFunction(s);
    }

    // 得到字符串 s 的每个后缀与 t 的前缀匹配长度
    std::vector<int> getMatch(const std::string &s, const std::string &t) {
        auto z = zFunction(t + "#" + s);
        return std::vector<int>(z.begin() + t.size() + 1, z.end());
    }

private:
    std::vector<int> zFunction(const std::string &s) {
        int n = s.size();
        std::vector<int> z(n);
        int l = 0, r = 0;
        for (int i = 1; i < n; ++i) {
            if (i < r) {
                z[i] = std::min(z[i - l], r - i);
            }
            while (i + z[i] < n && s[i + z[i]] == s[z[i]]) {
                ++z[i];
            }
            if (i + z[i] > r) {
                l = i;
                r = i + z[i];
            }
        }
        return z;
    }
};
    \end{lstlisting}

\subsection{Manacher}

    \txt{Manacher 算法在线性时间内求出每个位置为中心的最长回文半径。常见做法是在原串间插入分隔符，统一奇偶回文的处理。}

    \begin{lstlisting}
std::vector<int> manacher(const std::string &s) {
    std::string t;
    t.reserve(s.size() * 2 + 1);
    t.push_back('#');
    for (char c : s) {
        t.push_back(c);
        t.push_back('#');
    }
    int n = t.size();
    std::vector<int> r(n);
    int center = 0;
    for (int i = 0; i < n; ++i) {
        if (2 * center - i >= 0 && center + r[center] > i) {
            r[i] = std::min(r[2 * center - i], center + r[center] - i);
        }
        while (i - r[i] >= 0 && i + r[i] < n && t[i - r[i]] == t[i + r[i]]) {
            ++r[i];
        }
        if (i + r[i] > center + r[center]) {
            center = i;
        }
    }
    for (auto &x : r) --x;  // 去掉分隔符的贡献
    return r;
}
    \end{lstlisting}

\subsection{Minimal period | 最小循环节}

    \txt{判断一个串是否有长度为 $d$ 的循环节，只需比较区间 $[l,r-d]$ 与 $[l+d,r]$ 是否完全相同。}
    \txt{最小循环节长度一定是 $len$ 的因子，可用字符串哈希配合质因数分解不断尝试缩小周期。}

    \begin{lstlisting}
int findMinimalCycle(std::string s) {
    int len = (int)s.length();
    s = ' ' + s;  // 1-index
    int l = 1, r = len;
    StringHash hash(s, 13331, 998244353);
    int ans = len;
    int tmp = len;
    while (tmp != 1) {
        int p = minPrimeFactor(tmp);
        int d = ans / p;
        if (hash(l, r - d) == hash(l + d, r)) {
            ans = d;
        }
        tmp /= p;
    }
    return ans;
}
    \end{lstlisting}

\subsection{Lyndon decomposition | Lyndon 分解}

    \txt{Lyndon 串定义为本身严格小于其所有非空真后缀的串。任意串都存在唯一的 Lyndon 分解
    \[
        s = w_1 w_2 \dots w_k,\quad w_1 \ge w_2 \ge \dots \ge w_k,
    \]
    其中每个 $w_i$ 是 Lyndon 串。}

    \begin{lstlisting}
// 返回每个 Lyndon 块的右端点（1-index，下标相对原串）
std::vector<int> Lyndon(std::string s) {
    std::vector<int> res;
    int n = s.length();
    s = ' ' + s;
    for (int i = 1; i <= n; ) {
        int j = i, k = i + 1;
        while (k <= n && s[j] <= s[k]) {
            if (s[j] < s[k]) {
                j = i;
            } else {
                ++j;
            }
            ++k;
        }
        while (i <= j) {
            res.push_back(i + k - j - 1);
            i += k - j;
        }
    }
    return res;
}
    \end{lstlisting}

\subsection{Subsequence automaton | 子序列自动机}

    \txt{子序列自动机预处理 $next[i][c]$：从位置 $i$ 开始往后，第一个字符 $c$ 出现的位置。朴素做法时间复杂度 $O(n\cdot \Sigma)$。}

    \begin{lstlisting}
template <int Z, char Base>
struct SubSequenceAutomaton {
    std::vector<std::array<int, Z>> next;
    int n;

    // s 默认加过前缀空字符，1-index
    explicit SubSequenceAutomaton(const std::string &s)
        : n((int)s.size() - 1), next(n + 1) {
        for (int c = 0; c < Z; ++c) {
            next[n][c] = n + 1;
        }
        for (int i = n - 1; i >= 0; --i) {
            next[i] = next[i + 1];
            next[i][s[i + 1] - Base] = i + 1;
        }
    }

    int getNextPos(int pos, char c) const {
        return next[pos][c - Base];
    }
};
    \end{lstlisting}

    \txt{当字符集很大时，可以用可持久化线段树压缩值域，仅在发生改变的位置更新，查询复杂度降为 $O(\log n)$。}

    \subsubsection{Subsequence automaton with President tree | 主席树优化子序列自动机}

        \begin{lstlisting}
struct PresidentTree {
    int idx = 0;
    int n;
    std::vector<int> root, lc, rc;
    std::vector<int> sum;

    void pushup(int u) {
        sum[u] = sum[lc[u]] + sum[rc[u]];
    }

    // a 为值域数组，len 为时间轴长度
    PresidentTree(const std::vector<int> &a, int len)
        : n(a.size()), root(len + 1),
          lc(len * 40), rc(len * 40), sum(len * 40) {
        // 传入的 a 实际上是值域
        std::function<void(int &, int, int)> build =
            [&](int &u, int l, int r) {
                u = ++idx;
                if (l == r) {
                    sum[u] = a[l];
                    return;
                }
                int mid = (l + r) >> 1;
                build(lc[u], l, mid);
                build(rc[u], mid + 1, r);
                pushup(u);
            };
        build(root[len], 1, n - 1);  // 倒着建树
    }

    void insert(int &u, int v, int l, int r, int pos, int x) {
        u = ++idx;
        lc[u] = lc[v];
        rc[u] = rc[v];
        sum[u] = sum[v];
        if (l == r) {
            sum[u] = x;
            return;
        }
        int mid = (l + r) >> 1;
        if (pos <= mid) {
            insert(lc[u], lc[v], l, mid, pos, x);
        } else {
            insert(rc[u], rc[v], mid + 1, r, pos, x);
        }
        pushup(u);
    }

    void insert(int now, int pre, int pos, int x) {
        insert(root[now], root[pre], 1, n - 1, pos, x);
    }

    int query(int u, int l, int r, int pos) {
        if (l == r) {
            return sum[u];
        }
        int mid = (l + r) >> 1;
        if (pos <= mid) {
            return query(lc[u], l, mid, pos);
        } else {
            return query(rc[u], mid + 1, r, pos);
        }
    }

    int query(int t, int pos) {
        return query(root[t], 1, n - 1, pos);
    }
};

// 子序列自动机（值域大时）
struct SubsequenceAutomaton {
    const std::vector<int> &val;
    PresidentTree Tree;
    int n;

    SubsequenceAutomaton(const std::vector<int> &a, std::vector<int> &init)
        : val(a), Tree(init, (int)a.size() - 1),
          n((int)a.size() - 1) {
        for (int i = n - 1; i >= 0; --i) {
            // 更新子序列自动机
            Tree.insert(i, i + 1, val[i + 1], i + 1);
        }
    }

    // pos 后面第一个值为 val 的位置
    int query(int pos, int v) {
        return Tree.query(pos, v);
    }
};
        \end{lstlisting}

\subsection{AC automaton | AC 自动机}

    \txt{AC 自动机是在 \texttt{trie} 树基础上加上 \emph{fail} 指针得到的有向图。}
    \txt{对于节点 $p$，若沿字符 $c$ 没有儿子，则跳到 $fail(p)$ 再尝试 $c$；fail 树上父亲表示更短的后缀。多模式匹配时在 AC 图上跑一遍文本串，并把贡献沿 fail 树向上累加即可。}

    \begin{lstlisting}
template <int Z, char Base>
struct AcAutomaton {
    std::vector<int> fail;                 // fail 指针
    std::vector<std::vector<int>> ch;      // trie 图
    std::vector<std::vector<int>> ID;      // 终止串编号
    int idx = 0;
    int sum = 0;
    int n;                                 // 模式串个数

    explicit AcAutomaton(std::vector<std::string> &s)
        : n((int)s.size()) {
        for (auto &x : s) {
            sum += (int)x.size();
        }
        ch.assign(sum + 1, std::vector<int>(Z, 0));
        fail.assign(sum + 1, 0);
        ID.resize(sum + 1);
        for (int i = 0; i < n; ++i) {
            insert(i, s[i]);
        }
        build();
    }

    void insert(int id, const std::string &s) {
        int p = 0;
        for (char c : s) {
            int x = c - Base;
            if (!ch[p][x]) {
                ch[p][x] = ++idx;
            }
            p = ch[p][x];
        }
        ID[p].push_back(id);
    }

    void build() {
        std::queue<int> q;
        for (int i = 0; i < Z; ++i) {
            if (ch[0][i]) {
                q.push(ch[0][i]);
            }
        }
        while (!q.empty()) {
            int u = q.front();
            q.pop();
            for (int i = 0; i < Z; ++i) {
                int v = ch[u][i];
                if (v) {
                    fail[v] = ch[fail[u]][i];
                    q.push(v);
                } else {
                    ch[u][i] = ch[fail[u]][i];
                }
            }
        }
    }
};
    \end{lstlisting}

\subsection{Suffix automaton | 后缀自动机}

    \txt{后缀自动机（SAM）是一个 DAG，加上一棵 \emph{fail} 树。}
    \txt{每个状态代表一个 endpos 等价类，\texttt{len[u]} 为该状态表示的最长子串长度，\texttt{fail[u]} 指向最长真后缀所在的状态。}

    \begin{lstlisting}
template <int Z, char Base>
struct Sam {
    std::vector<std::array<int, Z>> ch;  // 转移
    std::vector<std::vector<int>> g;     // fail 树
    std::vector<int> fail, len;
    std::vector<long long> cnt;          // endpos 大小

    int last;      // 上一个添加的状态
    int tot;       // 总状态数
    long long num; // 不同子串数量

    explicit Sam(int n)
        : ch((n + 1) << 1),
          g((n + 1) << 1),
          fail((n + 1) << 1),
          len((n + 1) << 1),
          cnt((n + 1) << 1, 0),
          last(1), tot(1), num(0) {
        fail[1] = -1;
    }

    explicit Sam(const std::string &s)
        : Sam((int)s.size()) {
        for (char c : s) {
            add(c);
        }
        for (int i = 2; i <= tot; ++i) {
            g[fail[i]].push_back(i);
        }
    }

    void add(char c) {
        c -= Base;
        int cur = ++tot;
        len[cur] = len[last] + 1;
        cnt[cur] = 1;
        int v = last;
        while (v != -1 && !ch[v][c]) {
            ch[v][c] = cur;
            v = fail[v];
        }
        if (v == -1) {
            fail[cur] = 1;
        } else {
            int q = ch[v][c];
            if (len[v] + 1 == len[q]) {
                fail[cur] = q;
            } else {
                int clone = ++tot;
                len[clone] = len[v] + 1;
                ch[clone] = ch[q];
                fail[clone] = fail[q];
                while (v != -1 && ch[v][c] == q) {
                    ch[v][c] = clone;
                    v = fail[v];
                }
                fail[q] = fail[cur] = clone;
            }
        }
        last = cur;
        num += len[cur] - (fail[cur] == -1 ? 0 : len[fail[cur]]);
    }
};
    \end{lstlisting}

    \subsubsection{Dynamic SAM with LCT (sketch) | LCT 维护 SAM 链 / 子树信息}

    \txt{当需要在线维护串的改变时，可以用 LCT 维护 fail 树上的链或子树信息。核心是在 \texttt{add} 中对克隆节点做 \texttt{splay} 并复制节点信息，保证历史贡献不丢失。此处省略 LCT 具体实现，只保留 \texttt{add} 思路。}

\subsection{Generalized SAM | 广义后缀自动机}

    \txt{广义 SAM 在一棵 trie 上建多个字符串的 SAM，每个状态维护一个长度为串数的数组 \texttt{ID[u][i]}，表示该状态代表的子串在第 $i$ 个串中出现次数。}
    \txt{在 fail 树上自底向上累加可以得到每个状态在每个串中的 endpos 大小，从而解决多串最长公共子串等问题。}

    \begin{lstlisting}
template <int Z, char Base>
struct ExSam {
    std::vector<std::array<int, Z>> ch;
    std::vector<int> len, fail;
    std::vector<std::vector<int>> ID;
    std::vector<std::vector<int>> g;
    int cnt;            // 字符串数量

    int tot = 0, last = 0;
    long long num = 0;

    explicit ExSam(const std::vector<std::string> &S) {
        cnt = (int)S.size();
        int LEN = 0;
        for (auto &str : S) {
            LEN += (int)str.size();
        }
        len.assign(LEN << 1, 0);
        fail.assign(LEN << 1, 0);
        ch.assign(LEN << 1, {});
        g.assign(LEN << 1, {});
        ID.assign(LEN << 1, std::vector<int>(cnt, 0));
        fail[0] = -1;

        for (int i = 0; i < cnt; ++i) {
            last = 0;
            for (char c : S[i]) {
                exadd(c, i);
            }
        }
    }

    void dfs(int u) {
        for (int v : g[u]) {
            dfs(v);
            for (int i = 0; i < cnt; ++i) {
                ID[u][i] += ID[v][i];
            }
        }
    }

    bool check(int p) {
        for (int i = 0; i < cnt; ++i) {
            if (!ID[p][i]) return false;
        }
        return true;
    }

    int querylen(const std::string &s) {
        for (int i = 1; i <= tot; ++i) {
            g[fail[i]].push_back(i);
        }
        dfs(0);
        int p = 0, now = 0, res = 0;
        for (char c : s) {
            int x = c - Base;
            if (check(ch[p][x])) {
                p = ch[p][x];
                ++now;
            } else {
                while (p != -1 && !check(ch[p][x])) {
                    p = fail[p];
                }
                if (p == -1) {
                    now = 0;
                    p = 0;
                } else {
                    now = len[p] + 1;
                    p = ch[p][x];
                }
            }
            res = std::max(res, now);
        }
        return res;
    }

private:
    void assign(int cur, int id) {
        ID[cur][id] += 1;
    }

    void exadd(char c, int id) {
        int x = c - Base;
        if (ch[last][x]) {
            int v = last;
            int q = ch[v][x];
            if (len[q] != len[v] + 1) {
                int cl = ++tot;
                len[cl] = len[v] + 1;
                fail[cl] = fail[q];
                ch[cl] = ch[q];
                while (v != -1 && ch[v][x] == q) {
                    ch[v][x] = cl;
                    v = fail[v];
                }
                fail[q] = cl;
                q = cl;
            }
            int cur = last = q;
            assign(cur, id);
            return;
        }
        int cur = ++tot;
        len[cur] = len[last] + 1;
        assign(cur, id);
        int v = last;
        while (v != -1 && !ch[v][x]) {
            ch[v][x] = cur;
            v = fail[v];
        }
        if (v == -1) {
            fail[cur] = 0;
        } else {
            int q = ch[v][x];
            if (len[q] == len[v] + 1) {
                fail[cur] = q;
            } else {
                int cl = ++tot;
                len[cl] = len[v] + 1;
                fail[cl] = fail[q];
                ch[cl] = ch[q];
                while (v != -1 && ch[v][x] == q) {
                    ch[v][x] = cl;
                    v = fail[v];
                }
                fail[q] = fail[cur] = cl;
            }
        }
        last = cur;
        num += len[cur] - len[fail[cur]];
    }
};
    \end{lstlisting}

\subsection{Palindromic automaton | 回文自动机}

    \txt{回文自动机（PAM）维护所有不同回文子串。建两棵根：偶数根 $0$ 和奇数根 $1$，其中 $len[0]=0,len[1]=-1$，且 $fail[0]=1$。每个节点代表一个回文串，\texttt{fail[u]} 为其最长真回文后缀。}

    \begin{lstlisting}
template <int Z, char Base>
struct Pam {
    std::vector<std::array<int, Z>> ch;
    std::vector<int> fail, len, dep, cnt;
    std::string s;

    int cur = 0, tot = 1;  // 0: 偶根, 1: 奇根

    explicit Pam(int n)
        : ch(n + 2), fail(n + 2),
          len(n + 2), dep(n + 2), cnt(n + 2, 0) {
        fail[0] = 1;
        len[1] = -1;
    }

    explicit Pam(const std::string &str)
        : Pam((int)str.size()) {
        for (int i = 0; i < (int)str.size(); ++i) {
            add(i, str[i]);
        }
    }

    void assign(int u, int /*id*/) {
        cnt[u] += 1;
    }

    int getfail(int x, int i) {
        while (i - len[x] - 1 < 0 || s[i - len[x] - 1] != s[i]) {
            x = fail[x];
        }
        return x;
    }

    void add(int i, char c) {
        int x = c - Base;
        s.push_back(c);
        int v = getfail(cur, i);
        if (!ch[v][x]) {
            fail[++tot] = ch[getfail(fail[v], i)][x];
            ch[v][x] = tot;
            len[tot] = len[v] + 2;
            dep[tot] = dep[fail[tot]] + 1;
        }
        cur = ch[v][x];
        assign(cur, i);
    }

    auto getFailTree() const {
        std::vector<std::vector<int>> e(tot + 1);
        for (int i = 0; i <= tot; ++i) {
            if (i != 1) {
                e[fail[i]].push_back(i);
            }
        }
        return e;
    }
};
    \end{lstlisting}

\subsection{Suffix array | 后缀数组}

    \txt{后缀数组维护后缀的字典序排序，\texttt{sa[i]} 为第 $i$ 小后缀的起始下标，\texttt{rk[i]} 为后缀 $i$ 的排名，\texttt{h[i]} 为 $sa[i]$ 与 $sa[i-1]$ 的 LCP。常用性质包括：}
    \txt{1. $LCP(sa[i],sa[j]) = \min\limits_{k=i+1}^{j} h[k]$。\quad 2. 不同子串数：$n(n+1)/2 - \sum_{i=1}^{n} h[i]$。}

    \begin{lstlisting}
struct SuffixArray {
    const int n;
    std::vector<int> sa, rk, h;

    template <class T>
    explicit SuffixArray(const T &s)
        : n((int)s.size()),
          sa(n), rk(n),
          id(n), tmp(n) {
        std::iota(id.begin(), id.end(), 0);
        for (int i = 0; i < n; ++i) {
            rk[i] = s[i];
        }
        countSort();
        for (int w = 1;; w <<= 1) {
            std::iota(id.begin(), id.begin() + w + 1, n - w);
            for (int i = 0, p = w; i < n; ++i) {
                if (sa[i] >= w) {
                    id[p++] = sa[i] - w;
                }
            }
            countSort();
            oldrk = rk;
            rk[sa[0]] = 0;
            for (int i = 1, p = 0; i < n; ++i) {
                rk[sa[i]] = equal(sa[i], sa[i - 1], w) ? p : ++p;
            }
            if (rk[sa.back()] + 1 == n) break;
        }
        calcHeight(s);
    }

private:
    std::vector<int> oldrk, id, tmp, cnt;

    template <class T>
    void calcHeight(const T &s) {
        h.assign(n, 0);
        for (int i = 0, k = 0; i < n; ++i) {
            if (!rk[i]) continue;
            if (k) --k;
            while (i + k < n && s[i + k] == s[sa[rk[i] - 1] + k]) {
                ++k;
            }
            h[rk[i]] = k;
        }
    }

    void countSort() {
        int m = *std::max_element(rk.begin(), rk.end());
        cnt.assign(m + 1, 0);
        for (int i = 0; i < n; ++i) {
            cnt[tmp[i] = rk[id[i]]] += 1;
        }
        for (int i = 1; i < (int)cnt.size(); ++i) {
            cnt[i] += cnt[i - 1];
        }
        for (int i = n - 1; i >= 0; --i) {
            sa[--cnt[tmp[i]]] = id[i];
        }
    }

    bool equal(int x, int y, int w) {
        int rkx = (x + w < n ? oldrk[x + w] : -1);
        int rky = (y + w < n ? oldrk[y + w] : -1);
        return oldrk[x] == oldrk[y] && rkx == rky;
    }
};
    \end{lstlisting}



\section{Number Theory | 数论}

    \subsection{ModInt | 模运算}
        \code{src/number_theory/modint.cpp}

    \subsection{Prime sieve | 素数筛}
        \code{src/number_theory/prime_sieve.cpp}

    \subsection{Qpow block | 快速幂分块}
        \code{src/number_theory/qpow_block.cpp}

    \subsection{O(1) GCD | O(1) 最大公约数}
            \txt{适用于\inlineformula{./assets/formulas_tex/formula_001.tex}以内的数字求\inlineformula{./assets/formulas_tex/formula_002.tex},我们传入的\inlineformula{./assets/formulas_tex/formula_003.tex}为能接受数字的上限，\inlineformula{./assets/formulas_tex/formula_004.tex}为小表参数，复杂度为\inlineformula{./assets/formulas_tex/formula_005.tex}，查询可以认为是\inlineformula{./assets/formulas_tex/formula_006.tex}，我们一般传入\inlineformula{./assets/formulas_tex/formula_007.tex}即可。}
            \code{src/number_theory/O(1)gcd_5.cpp}

    \subsection{Number theory basics | 数论基础知识}
        \subsubsection{关于狄利克雷卷积的定义}
        \subsubsection{质数表}
            \code{src/number_theory/数论基础知识_5.cpp}

    \subsection{Number theory theorems | 数论定理}
        \subsubsection{三平方和定理}
            \txt{对于任意自然数\inlineformula{./assets/formulas_tex/formula_008.tex}，可以表示为\displayformula{./assets/formulas_tex/formula_009.tex}形式的\displayformula{./assets/formulas_tex/formula_010.tex}的条件是\displayformula{./assets/formulas_tex/formula_011.tex}}
        \subsubsection{拉格朗日四平方和定理}
            \txt{对于任意自然数\displayformula{./assets/formulas_tex/formula_012.tex}，都可以表示为\displayformula{./assets/formulas_tex/formula_013.tex}，其中\displayformula{./assets/formulas_tex/formula_014.tex}可以为\displayformula{./assets/formulas_tex/formula_015.tex}
暴力代码
大约为\displayformula{./assets/formulas_tex/formula_016.tex}
优化代码
大约为\displayformula{./assets/formulas_tex/formula_017.tex}到\displayformula{./assets/formulas_tex/formula_018.tex}}
            \code{src/number_theory/数论定理_21.cpp}
        \subsubsection{雅可比四平方和定理}
            \txt{设能满足\displayformula{./assets/formulas_tex/formula_019.tex}的自然数解个数为\displayformula{./assets/formulas_tex/formula_020.tex},那么如果\displayformula{./assets/formulas_tex/formula_021.tex}为奇数，\displayformula{./assets/formulas_tex/formula_022.tex}，否则我们记\displayformula{./assets/formulas_tex/formula_023.tex},即去掉所有\displayformula{./assets/formulas_tex/formula_024.tex}的质因子，此时\displayformula{./assets/formulas_tex/formula_025.tex}。其中\displayformula{./assets/formulas_tex/formula_026.tex}表示\displayformula{./assets/formulas_tex/formula_027.tex}的约数和}
        \subsubsection{毕达哥拉斯三元组}
            \txt{\displayformula{./assets/formulas_tex/formula_028.tex}
令\displayformula{./assets/formulas_tex/formula_029.tex}
如果要求\displayformula{./assets/formulas_tex/formula_030.tex},条件是\displayformula{./assets/formulas_tex/formula_031.tex}和\displayformula{./assets/formulas_tex/formula_032.tex}至少一个是偶数且\displayformula{./assets/formulas_tex/formula_033.tex}}
        \subsubsection{预处理因数}
            \txt{我们可以在\inlineformula{./assets/formulas_tex/formula_034.tex}的时间内处理出\inlineformula{./assets/formulas_tex/formula_035.tex}中所有数字的所有因数}
            \code{src/number_theory/数论定理_3_5.cpp}
        \subsubsection{欧拉筛}
            \code{src/number_theory/数论定理_4_5.cpp}
        \subsubsection{欧拉函数}
            \txt{欧拉函数的数论意义: \displayformula{./assets/formulas_tex/formula_036.tex}表示在\displayformula{./assets/formulas_tex/formula_037.tex}中和\inlineformula{./assets/formulas_tex/formula_038.tex}互质的数字的数量
公式：\inlineformula{./assets/formulas_tex/formula_039.tex} 其中\inlineformula{./assets/formulas_tex/formula_040.tex}为x的质因子 \inlineformula{./assets/formulas_tex/formula_041.tex}为质因子数量}
            \code{src/number_theory/数论定理_5_5.cpp}
        \subsubsection{莫比乌斯函数}
            \txt{莫比乌斯函数最重要的性质：\displayformula{./assets/formulas_tex/formula_042.tex} 用莫比乌斯反演来写 也就是 \displayformula{./assets/formulas_tex/formula_043.tex}
\displayformula{./assets/formulas_tex/formula_044.tex}}
            \code{src/number_theory/数论定理_6_5.cpp}
        \subsubsection{欧拉定理}
        \subsubsection{扩展欧拉定理}
        \subsubsection{扩展欧几里得算法}
            \txt{即求 \displayformula{./assets/formulas_tex/formula_045.tex}的通解  根据裴蜀定理我们知道 当\inlineformula{./assets/formulas_tex/formula_046.tex}
所以该二元不定方程有解当且仅当\inlineformula{./assets/formulas_tex/formula_047.tex} 我们考虑求出简单式的解再乘\inlineformula{./assets/formulas_tex/formula_048.tex}即可
同时我们发现 假设一组解\inlineformula{./assets/formulas_tex/formula_049.tex}是方程\inlineformula{./assets/formulas_tex/formula_050.tex}的一组特解
那么我们可以写出其通解  通解形式为\inlineformula{./assets/formulas_tex/formula_051.tex}
那么我们考虑求出一组特解即可
\inlineformula{./assets/formulas_tex/formula_052.tex}
再次根据裴蜀定理我们有
\inlineformula{./assets/formulas_tex/formula_053.tex}
结合两个式子我们有 \inlineformula{./assets/formulas_tex/formula_054.tex}  那么根据欧几里得算法的过程一定会递归到\inlineformula{./assets/formulas_tex/formula_055.tex}
那么此时结果为\inlineformula{./assets/formulas_tex/formula_056.tex} 至此我们的特解求解完成 最终通解仍然要乘\inlineformula{./assets/formulas_tex/formula_057.tex}}
            \code{src/number_theory/数论定理_7_5.cpp}
        \subsubsection{扩展中国剩余定理}
            \txt{扩展中国剩余定理 是用来求解一次同余方程组的通解 考虑到\inlineformula{./assets/formulas_tex/formula_058.tex}会比\inlineformula{./assets/formulas_tex/formula_059.tex}适用性更广
（\inlineformula{./assets/formulas_tex/formula_060.tex}数可以不互质） 我们直接考虑\inlineformula{./assets/formulas_tex/formula_061.tex}
方程形式即 \inlineformula{./assets/formulas_tex/formula_062.tex}(我们用\inlineformula{./assets/formulas_tex/formula_063.tex}表示都满足)
我么考虑取出前两个式子 \inlineformula{./assets/formulas_tex/formula_064.tex}，\displayformula{./assets/formulas_tex/formula_065.tex}
\displayformula{./assets/formulas_tex/formula_066.tex}
我们做差 得到\inlineformula{./assets/formulas_tex/formula_067.tex} 我们通过Exgcd解出 \inlineformula{./assets/formulas_tex/formula_068.tex}和\inlineformula{./assets/formulas_tex/formula_069.tex}的通解
我们得到\inlineformula{./assets/formulas_tex/formula_070.tex}和\inlineformula{./assets/formulas_tex/formula_071.tex}后再次代入\inlineformula{./assets/formulas_tex/formula_072.tex}中 得到
\inlineformula{./assets/formulas_tex/formula_073.tex}
其中\displayformula{./assets/formulas_tex/formula_074.tex}表示刚刚解出的\displayformula{./assets/formulas_tex/formula_075.tex}的特解
不难发现我们再次得到了一个同余式\displayformula{./assets/formulas_tex/formula_076.tex}
我们不断合并直至只剩一个式子 即得答案 复杂度\displayformula{./assets/formulas_tex/formula_077.tex}}
            \code{src/number_theory/数论定理_8_5.cpp}
        \subsubsection{扩展大步小步算法}
            \txt{我们先提出BSGS算法
BSGS的算法思想很多时候能发挥重要作用，其中文名为大步小步算法，实际上是在进行一种优雅的暴力枚举
我们的目的是求 \displayformula{./assets/formulas_tex/formula_078.tex}的解\inlineformula{./assets/formulas_tex/formula_079.tex}
我们设\inlineformula{./assets/formulas_tex/formula_080.tex} 其中\inlineformula{./assets/formulas_tex/formula_081.tex}我们称这个m为大步
我们移项以后\inlineformula{./assets/formulas_tex/formula_082.tex}我们称1为小步
我们枚举\inlineformula{./assets/formulas_tex/formula_083.tex}从\inlineformula{./assets/formulas_tex/formula_084.tex}到\inlineformula{./assets/formulas_tex/formula_085.tex}，用哈希表存下所有左边的值，再枚举\inlineformula{./assets/formulas_tex/formula_086.tex}同样从\displayformula{./assets/formulas_tex/formula_087.tex}到\inlineformula{./assets/formulas_tex/formula_088.tex}，求出右边的值，在哈希表中查找是否存在答案即可，我们发现所有小步的长度和等于一大步，我们发现这覆盖的范围是\inlineformula{./assets/formulas_tex/formula_089.tex}，当\inlineformula{./assets/formulas_tex/formula_090.tex}，由于扩展欧拉定理的内容,我们发现这已经可以覆盖所有情况，复杂度为\inlineformula{./assets/formulas_tex/formula_091.tex}
然后我们考虑如何解决(a,P)不互质的情况得到EXBSGS算法
考虑把同余式转化为方程式，即\inlineformula{./assets/formulas_tex/formula_092.tex}我们设\inlineformula{./assets/formulas_tex/formula_093.tex},那么由裴蜀定理
当\inlineformula{./assets/formulas_tex/formula_094.tex}的时候原方程有解,如果有解，我们可以转化为 \displayformula{./assets/formulas_tex/formula_095.tex} 我们重复此过程直至
\inlineformula{./assets/formulas_tex/formula_096.tex}  我们设递归了\inlineformula{./assets/formulas_tex/formula_097.tex}次，这\displayformula{./assets/formulas_tex/formula_098.tex}次\displayformula{./assets/formulas_tex/formula_099.tex}的乘积为\inlineformula{./assets/formulas_tex/formula_100.tex}
那么原方程为\inlineformula{./assets/formulas_tex/formula_101.tex} 注意到此时变成了一个BSGS问题但前面多了一个常数
我们可以求其逆元丢到右边即可}
            \code{src/number_theory/数论定理_9_5.cpp}
        \subsubsection{卢卡斯定理}
            \txt{即 \inlineformula{./assets/formulas_tex/formula_102.tex}
不难发现我们写个递归函数就行，请注意递归出口是\inlineformula{./assets/formulas_tex/formula_103.tex}，此时返回\inlineformula{./assets/formulas_tex/formula_104.tex}。}
        \subsubsection{原根}
            \txt{阶的定义 ：在\inlineformula{./assets/formulas_tex/formula_105.tex}意义下 满足\inlineformula{./assets/formulas_tex/formula_106.tex}的最小正整数\inlineformula{./assets/formulas_tex/formula_107.tex}称作\inlineformula{./assets/formulas_tex/formula_108.tex}\inlineformula{./assets/formulas_tex/formula_109.tex} 意义下的阶
原根的定义：一个数在\inlineformula{./assets/formulas_tex/formula_110.tex}意义下的阶等于 \inlineformula{./assets/formulas_tex/formula_111.tex} 那么该数称作\inlineformula{./assets/formulas_tex/formula_112.tex}的一个原根
首先已有数学家证明 一个数\inlineformula{./assets/formulas_tex/formula_113.tex}的最小原根不会超过\inlineformula{./assets/formulas_tex/formula_114.tex}
所以我们考虑从小到大暴力枚举
在此之前我们已经证明 只有 \inlineformula{./assets/formulas_tex/formula_115.tex}存在原根（这里P是质数）
所以我们先通过欧拉筛预处理 判断出一个数是否存在原根
在存在的情况下 我们先求出最小原根  从1开始枚举 然后对每个数分别Check
假设我们枚举到数字\inlineformula{./assets/formulas_tex/formula_116.tex} 我们枚举\inlineformula{./assets/formulas_tex/formula_117.tex}的所有质因数\inlineformula{./assets/formulas_tex/formula_118.tex} 如果\inlineformula{./assets/formulas_tex/formula_119.tex}也是\inlineformula{./assets/formulas_tex/formula_120.tex}的阶  那么\inlineformula{./assets/formulas_tex/formula_121.tex}就不是\inlineformula{./assets/formulas_tex/formula_122.tex}的原根
当我们找到最小原根\inlineformula{./assets/formulas_tex/formula_123.tex}后 对于所有的\inlineformula{./assets/formulas_tex/formula_124.tex}满足\inlineformula{./assets/formulas_tex/formula_125.tex}的所有\inlineformula{./assets/formulas_tex/formula_126.tex} ，满足\inlineformula{./assets/formulas_tex/formula_127.tex}也是原根
至此我们得到了求出所有原根的方法
且根据最后一个定理 我们知道 一个数\inlineformula{./assets/formulas_tex/formula_128.tex}如果存在原根 那么它将有\inlineformula{./assets/formulas_tex/formula_129.tex}个原根}
            \code{src/number_theory/数论定理_10_5.cpp}
        \subsubsection{Pell 方程}
            \txt{即\inlineformula{./assets/formulas_tex/formula_131.tex}，求最小正整数解\inlineformula{./assets/formulas_tex/formula_132.tex}
由于答案可以很大，代码由\inlineformula{./assets/formulas_tex/formula_133.tex}实现
具体求解方法涉及连分数，这里不过多介绍}
            \code{src/number_theory/数论定理_11_5.cpp}
        \subsubsection{类欧几里得算法}
            \txt{类欧算法是一类特定式子的求和方法
\inlineformula{./assets/formulas_tex/formula_134.tex}
\inlineformula{./assets/formulas_tex/formula_135.tex}
\inlineformula{./assets/formulas_tex/formula_136.tex}
注意代码并没有处理取模的问题
第一段代码只解决了\inlineformula{./assets/formulas_tex/formula_137.tex}的值
递推公式为
若\inlineformula{./assets/formulas_tex/formula_138.tex}
\inlineformula{./assets/formulas_tex/formula_139.tex}
否则
\inlineformula{./assets/formulas_tex/formula_140.tex}，其中\inlineformula{./assets/formulas_tex/formula_141.tex}
递归出口为\displayformula{./assets/formulas_tex/formula_142.tex}
第二段代码同时解决了\displayformula{./assets/formulas_tex/formula_143.tex}
推导公式过于复杂不再描述}
            \code{src/number_theory/数论定理_12_5.cpp}
        \subsubsection{整除分块}
            \txt{我们考虑如何解决各式各样的整除分块
我们先考虑一维的情况
即\inlineformula{./assets/formulas_tex/formula_144.tex}其中m是常数 \displayformula{./assets/formulas_tex/formula_145.tex}为单元函数且单调增加,我们只需要每次求出右端点\inlineformula{./assets/formulas_tex/formula_146.tex}，然后\displayformula{./assets/formulas_tex/formula_147.tex}即可，问题在于如何求出\inlineformula{./assets/formulas_tex/formula_148.tex}，我们来推式子
首先我们已经知道左端点\inlineformula{./assets/formulas_tex/formula_149.tex}，也就相当于知道了当前的值\inlineformula{./assets/formulas_tex/formula_150.tex}也就是说\inlineformula{./assets/formulas_tex/formula_151.tex}那么\inlineformula{./assets/formulas_tex/formula_152.tex}应该满足\inlineformula{./assets/formulas_tex/formula_153.tex} 那么\inlineformula{./assets/formulas_tex/formula_154.tex}
然后我们考虑二维的情况，一般形式为\inlineformula{./assets/formulas_tex/formula_155.tex}，在这种情况下我们发现有两个限制，但我们只需要每次操作取他们的交集即可，因为有\inlineformula{./assets/formulas_tex/formula_156.tex}在。可以保证全覆盖
也就是取\inlineformula{./assets/formulas_tex/formula_157.tex}，至于函数的形式我们参考一维即可
上取整的整除分块}
            \code{src/number_theory/数论定理_14_5.cpp}
        \subsubsection{莫比乌斯反演}
            \txt{其实就是一个式子
如果\inlineformula{./assets/formulas_tex/formula_158.tex}那么有\inlineformula{./assets/formulas_tex/formula_159.tex}
当然也和\inlineformula{./assets/formulas_tex/formula_160.tex}等价}
        \subsubsection{杜教筛}
            \txt{杜教筛是一种能在亚线性时间内求出积性函数前缀和的方法
构造方法如下
我们设\inlineformula{./assets/formulas_tex/formula_161.tex}表示所求函数，\inlineformula{./assets/formulas_tex/formula_162.tex}表示辅助函数，\inlineformula{./assets/formulas_tex/formula_163.tex}表示前缀和函数
那么我们有如下式子
\inlineformula{./assets/formulas_tex/formula_164.tex}
加入我们可以构造恰当的数论函数\inlineformula{./assets/formulas_tex/formula_165.tex}使得:
1.可以快速计算\inlineformula{./assets/formulas_tex/formula_166.tex}
2可以快速计算\inlineformula{./assets/formulas_tex/formula_167.tex}的前缀和，从而用数论分块来递归求解\inlineformula{./assets/formulas_tex/formula_168.tex}
我们发现构造是一个困难的事情 我们在此给出欧拉函数和莫比乌斯函数的板子
对于欧拉函数 我们可以构造\inlineformula{./assets/formulas_tex/formula_169.tex}
对于莫比乌斯函数 我们可以构造 \inlineformula{./assets/formulas_tex/formula_170.tex}
套公式即可}
            \code{src/number_theory/数论定理_16_5.cpp}
        \subsubsection{Min25筛}
            \txt{min25筛是统计一类函数前缀和的方法
这类函数有如下特征
1 \inlineformula{./assets/formulas_tex/formula_171.tex}可以被表示成简单多项式的形式
2 \inlineformula{./assets/formulas_tex/formula_172.tex}需要容易求得
Min25筛的核心思想是把质数和合数的贡献分开统计，这样同时解决了一些函数质数前缀和的问题
我们设我们要求\displayformula{./assets/formulas_tex/formula_173.tex}
我们先考虑质数如何计算 设\inlineformula{./assets/formulas_tex/formula_174.tex}其中\inlineformula{./assets/formulas_tex/formula_175.tex}
共\inlineformula{./assets/formulas_tex/formula_176.tex}种取值
我们考虑DP,设\inlineformula{./assets/formulas_tex/formula_177.tex}其中\inlineformula{./assets/formulas_tex/formula_178.tex}表示\inlineformula{./assets/formulas_tex/formula_179.tex}的最小质因子，\inlineformula{./assets/formulas_tex/formula_180.tex}是另一个我们构造出的函数，满足以下条件
\inlineformula{./assets/formulas_tex/formula_181.tex}
\inlineformula{./assets/formulas_tex/formula_182.tex}是完全积性函数
\inlineformula{./assets/formulas_tex/formula_183.tex}可以快速求前缀和
DP边界为\inlineformula{./assets/formulas_tex/formula_184.tex}
那么我们的转移有
\inlineformula{./assets/formulas_tex/formula_185.tex}因为这里显然不会多出数字满足\inlineformula{./assets/formulas_tex/formula_186.tex}
\inlineformula{./assets/formulas_tex/formula_187.tex}
这里的意义是减去最小质因子为\inlineformula{./assets/formulas_tex/formula_188.tex}的数字的贡献 利用了完全积性函数的性质
令\inlineformula{./assets/formulas_tex/formula_189.tex}表示小于等于\inlineformula{./assets/formulas_tex/formula_190.tex}的质数个数,那么\inlineformula{./assets/formulas_tex/formula_191.tex}即为我们所求的\inlineformula{./assets/formulas_tex/formula_192.tex}
当然我们会把第一维滚掉 用整除分块求出所有\inlineformula{./assets/formulas_tex/formula_193.tex}并用\inlineformula{./assets/formulas_tex/formula_194.tex}数组映射
至此我们已经完成了质数前缀和的工作，我们想求出所有数的前缀和，还需要二次\inlineformula{./assets/formulas_tex/formula_195.tex}
我们设\inlineformula{./assets/formulas_tex/formula_196.tex}那么最终答案为\inlineformula{./assets/formulas_tex/formula_197.tex}
有递推式
\inlineformula{./assets/formulas_tex/formula_198.tex}
请注意 对\inlineformula{./assets/formulas_tex/formula_199.tex}求出的\inlineformula{./assets/formulas_tex/formula_200.tex}筛只能得到所有的\inlineformula{./assets/formulas_tex/formula_201.tex}位置的质数前缀和
在模板中 \inlineformula{./assets/formulas_tex/formula_202.tex}，对于不同的函数我们只需要修改\inlineformula{./assets/formulas_tex/formula_203.tex}函数
注意取模！！！}
            \code{src/number_theory/数论定理_17_5.cpp}
        \subsubsection{常见的积性函数}
            \txt{\inlineformula{./assets/formulas_tex/formula_204.tex}
\inlineformula{./assets/formulas_tex/formula_205.tex}
\inlineformula{./assets/formulas_tex/formula_206.tex}
\inlineformula{./assets/formulas_tex/formula_207.tex}
\inlineformula{./assets/formulas_tex/formula_208.tex}
\inlineformula{./assets/formulas_tex/formula_209.tex}
\inlineformula{./assets/formulas_tex/formula_210.tex}}

    \subsection{Combination preprocessing | 组合数预处理}
        \subsubsection{组合数预处理}
            \txt{我们考虑\inlineformula{./assets/formulas_tex/formula_211.tex}线性递推预处理出组合数，我们用\inlineformula{./assets/formulas_tex/formula_212.tex}表示阶乘，\inlineformula{./assets/formulas_tex/formula_213.tex}表示阶乘的逆元，\inlineformula{./assets/formulas_tex/formula_214.tex}表示逆元。那么有如下递推式：
\inlineformula{./assets/formulas_tex/formula_215.tex}
\inlineformula{./assets/formulas_tex/formula_216.tex}
\inlineformula{./assets/formulas_tex/formula_217.tex}
不难发现我们只要用快速幂求出最后一个\inlineformula{./assets/formulas_tex/formula_218.tex}的值即可。}
            \code{src/number_theory/组合数预处理_7.cpp}
        \subsubsection{取模类}
            \code{src/number_theory/组合数预处理_2_5.cpp}
        \subsubsection{枚举固定大小子集}
            \code{src/number_theory/组合数预处理_3_5.cpp}

    \subsection{Tricks | 技巧}
        \subsubsection{Miller Rabin}
            \code{src/number_theory/技巧_Trick_2_5.cpp}
        \subsubsection{Pollard Rho}
            \txt{如果 \inlineformula{./assets/formulas_tex/formula_219.tex} 是质数（Miller Rabin 判）返回 \inlineformula{./assets/formulas_tex/formula_220.tex}，否则随机返回 \inlineformula{./assets/formulas_tex/formula_221.tex} 的一个 \inlineformula{./assets/formulas_tex/formula_222.tex} 因子，复杂度 \inlineformula{./assets/formulas_tex/formula_223.tex}。}
            \code{src/number_theory/技巧_Trick_3_5.cpp}
        \subsubsection{随机数生成}
            \code{src/number_theory/技巧_Trick_4_5.cpp}

    \subsection{Classic counting models | 几个经典计数模型}
        \subsubsection{卡特兰数}
            \txt{我们考虑推导，首先我们考虑所有情况，那么显然为\displayformula{./assets/formulas_tex/formula_224.tex}，但是我们不能越过\inlineformula{./assets/formulas_tex/formula_225.tex}，也就是不能触碰\inlineformula{./assets/formulas_tex/formula_226.tex}，考虑求出不合法的情况，我们把触碰\inlineformula{./assets/formulas_tex/formula_227.tex}以后得折线部分关于\inlineformula{./assets/formulas_tex/formula_228.tex}对称，发现终点变成了\inlineformula{./assets/formulas_tex/formula_229.tex}，那么方案数为\inlineformula{./assets/formulas_tex/formula_230.tex}，总方案数为\inlineformula{./assets/formulas_tex/formula_231.tex}。
我们考虑拓展，假设有\inlineformula{./assets/formulas_tex/formula_232.tex}个\inlineformula{./assets/formulas_tex/formula_233.tex}，\inlineformula{./assets/formulas_tex/formula_234.tex}个\inlineformula{./assets/formulas_tex/formula_235.tex}，组成一个序列，要求前缀和恒大于等于\inlineformula{./assets/formulas_tex/formula_236.tex}，求方案数，发现总的情况数为\inlineformula{./assets/formulas_tex/formula_237.tex}，条件即不能触碰\inlineformula{./assets/formulas_tex/formula_238.tex}，那么我们考虑\inlineformula{./assets/formulas_tex/formula_239.tex}关于它的对称点，易得到为\inlineformula{./assets/formulas_tex/formula_240.tex}，那么方案数减去\inlineformula{./assets/formulas_tex/formula_241.tex}即可}
        \subsubsection{卡特兰数}
            \txt{常规的卡特兰数我们有两种表述形式，其一是\inlineformula{./assets/formulas_tex/formula_242.tex}式表述，其二是组合意义表述。
组合意义表述来说，就是存在两种元素，满足在序列的任意一个前缀中其中一种的个数恒大于等于另一种且最终二者数量相同 的方案数
就是常规卡特兰数了，例子很多，比如给一个进栈序列求出栈顺序的方案数，满足（入栈数>=出栈数 且最终相等），再比如括号正则表达式，满足(左括号数量>=右括号数量,且最终左等于右)
DP式表述来说，即形如\inlineformula{./assets/formulas_tex/formula_243.tex}的递推式，最终\inlineformula{./assets/formulas_tex/formula_244.tex}的表达式也将是卡特兰数，证明只需把通项公式展开即可,这个的例子形如n个点的二叉树有多少种形态，是由DP式子得到卡特兰数的结论
我们给出几个通项公式，\inlineformula{./assets/formulas_tex/formula_245.tex}
\inlineformula{./assets/formulas_tex/formula_246.tex}    \inlineformula{./assets/formulas_tex/formula_247.tex}}
        \subsubsection{卡特兰数数表}
            \code{src/number_theory/几个经典计数模型_6.cpp}
        \subsubsection{k阶卡特兰数}
            \txt{\inlineformula{./assets/formulas_tex/formula_248.tex}阶卡特兰数的定义是，在一个网格上我要走到\inlineformula{./assets/formulas_tex/formula_249.tex}点，我们移动的向量只有\inlineformula{./assets/formulas_tex/formula_250.tex},且要求我们不能越过\inlineformula{./assets/formulas_tex/formula_251.tex},求移动方案数。
显然当\inlineformula{./assets/formulas_tex/formula_252.tex}时方案数为常规卡特兰数，因为\inlineformula{./assets/formulas_tex/formula_253.tex}的数量恒大于等于\inlineformula{./assets/formulas_tex/formula_254.tex}且数量最终相等
当\inlineformula{./assets/formulas_tex/formula_255.tex}不再为\inlineformula{./assets/formulas_tex/formula_256.tex}时，问题变得复杂，我们直接给出公式和DP式，推导咕咕咕
\inlineformula{./assets/formulas_tex/formula_257.tex}
\inlineformula{./assets/formulas_tex/formula_258.tex}
可以解决如下问题
有\inlineformula{./assets/formulas_tex/formula_259.tex}个\inlineformula{./assets/formulas_tex/formula_260.tex},\inlineformula{./assets/formulas_tex/formula_261.tex}个\inlineformula{./assets/formulas_tex/formula_262.tex},要求你组成一个序列且该序列前缀和恒大于等于0，问序列个数}
        \subsubsection{欧拉错排数}
            \txt{即n个编号为\inlineformula{./assets/formulas_tex/formula_263.tex}的位置，我有同样编号为\displayformula{./assets/formulas_tex/formula_264.tex}的n个数，要求把n个数填到位置中，满足位置号和编号均不相同的方案数\displayformula{./assets/formulas_tex/formula_265.tex}
这个我们可以递推解决，满足\inlineformula{./assets/formulas_tex/formula_266.tex}
简单证明:我们假设第一个数放在位置\inlineformula{./assets/formulas_tex/formula_267.tex}上，\inlineformula{./assets/formulas_tex/formula_268.tex}有n-1种选择，我们考虑\inlineformula{./assets/formulas_tex/formula_269.tex}位置的选择 分两种情况
如果\inlineformula{./assets/formulas_tex/formula_270.tex}也放在\inlineformula{./assets/formulas_tex/formula_271.tex}位置上那么就是剩下\inlineformula{./assets/formulas_tex/formula_272.tex}个数再次错排，否则将是\inlineformula{./assets/formulas_tex/formula_273.tex}个数错排。}
        \subsubsection{莫茨金数 (Motzkin Numbers)}
            \txt{\inlineformula{./assets/formulas_tex/formula_274.tex}
\inlineformula{./assets/formulas_tex/formula_275.tex}
\inlineformula{./assets/formulas_tex/formula_276.tex}
经典应用场景：
\inlineformula{./assets/formulas_tex/formula_277.tex}一个圆上有\inlineformula{./assets/formulas_tex/formula_278.tex}个点，画若干条不相交弦的方案数
\inlineformula{./assets/formulas_tex/formula_279.tex}在一个二维网格从\inlineformula{./assets/formulas_tex/formula_280.tex}走到\displayformula{./assets/formulas_tex/formula_281.tex}，\displayformula{./assets/formulas_tex/formula_282.tex}表示向右上走一步，\displayformula{./assets/formulas_tex/formula_283.tex}表示向右走一步，\inlineformula{./assets/formulas_tex/formula_284.tex}表示向右下走一步
要求路径中始终不能低于\inlineformula{./assets/formulas_tex/formula_285.tex}轴，终点必须在\inlineformula{./assets/formulas_tex/formula_286.tex}，问方案数。}
        \subsubsection{贝尔数}
            \txt{贝尔数\inlineformula{./assets/formulas_tex/formula_287.tex}表示把\inlineformula{./assets/formulas_tex/formula_288.tex}个元素划分集合的方案数，集合数量无限制。
显然贝尔数是第二类斯特林数的和，即\inlineformula{./assets/formulas_tex/formula_289.tex}
递推公式\inlineformula{./assets/formulas_tex/formula_290.tex}
\inlineformula{./assets/formulas_tex/formula_291.tex}}
        \subsubsection{斐波那契数列}
        \subsubsection{斐波那契数表}
            \txt{斐波那契数列有许多有趣的性质，我们总结一下
\inlineformula{./assets/formulas_tex/formula_292.tex}斐波那契数列我们认为是三角形数，如果要构造一个最长的序列满足任意三个数都不能构成三角形的话，斐波那契数列将是最佳选择
\inlineformula{./assets/formulas_tex/formula_293.tex}\inlineformula{./assets/formulas_tex/formula_294.tex}
\inlineformula{./assets/formulas_tex/formula_295.tex}\inlineformula{./assets/formulas_tex/formula_296.tex}
\inlineformula{./assets/formulas_tex/formula_297.tex}\displayformula{./assets/formulas_tex/formula_298.tex} \inlineformula{./assets/formulas_tex/formula_299.tex}
\inlineformula{./assets/formulas_tex/formula_300.tex}\inlineformula{./assets/formulas_tex/formula_301.tex}
\inlineformula{./assets/formulas_tex/formula_302.tex}\inlineformula{./assets/formulas_tex/formula_303.tex}
\inlineformula{./assets/formulas_tex/formula_304.tex}\inlineformula{./assets/formulas_tex/formula_305.tex}
\inlineformula{./assets/formulas_tex/formula_306.tex}\inlineformula{./assets/formulas_tex/formula_307.tex}
\inlineformula{./assets/formulas_tex/formula_308.tex}\inlineformula{./assets/formulas_tex/formula_309.tex}
\inlineformula{./assets/formulas_tex/formula_310.tex}\inlineformula{./assets/formulas_tex/formula_311.tex}
\inlineformula{./assets/formulas_tex/formula_312.tex}斐波那契数列在\inlineformula{./assets/formulas_tex/formula_313.tex}意义在存在周期，称作皮萨诺周期，我们已经证明，\inlineformula{./assets/formulas_tex/formula_314.tex}\inlineformula{./assets/formulas_tex/formula_315.tex}的最小循环节不超过\inlineformula{./assets/formulas_tex/formula_316.tex}}
            \code{src/number_theory/几个经典计数模型_2_5.cpp}

    \subsection{Tree counting | 树计数}
        \subsubsection{无标号树计数}
            \txt{有根树计数：
\inlineformula{./assets/formulas_tex/formula_317.tex}
递推公式
\inlineformula{./assets/formulas_tex/formula_318.tex}
无根树计数：
\inlineformula{./assets/formulas_tex/formula_319.tex}为奇数时
\inlineformula{./assets/formulas_tex/formula_320.tex}
\inlineformula{./assets/formulas_tex/formula_321.tex}为偶数时
\inlineformula{./assets/formulas_tex/formula_322.tex}}
        \subsubsection{Prufer序列}
            \txt{\inlineformula{./assets/formulas_tex/formula_323.tex}序列实际上和无根树对应它是无根树衍生出来的一种序列
我们先给出转化方法
\inlineformula{./assets/formulas_tex/formula_324.tex}无根树转\inlineformula{./assets/formulas_tex/formula_325.tex}序列
首先我们找到编号最小且度数为\inlineformula{./assets/formulas_tex/formula_326.tex}的点
其次删除该节点并加入和该节点相连的父节点
我们重复上述操作直至剩下两个节点
不难发现\inlineformula{./assets/formulas_tex/formula_327.tex}个点的树的\inlineformula{./assets/formulas_tex/formula_328.tex}序列长度是\inlineformula{./assets/formulas_tex/formula_329.tex}的
因为每删一个点\inlineformula{./assets/formulas_tex/formula_330.tex}序列长度加一最后剩两个点说明了\inlineformula{./assets/formulas_tex/formula_331.tex}序列的长度为\inlineformula{./assets/formulas_tex/formula_332.tex};
\inlineformula{./assets/formulas_tex/formula_333.tex}\inlineformula{./assets/formulas_tex/formula_334.tex}序列转无根树
首先我们取出\inlineformula{./assets/formulas_tex/formula_335.tex}序列的首个元素 \inlineformula{./assets/formulas_tex/formula_336.tex}
其次在点集中找到最小的且未出现在\inlineformula{./assets/formulas_tex/formula_337.tex}序列中的点\inlineformula{./assets/formulas_tex/formula_338.tex} 连边\inlineformula{./assets/formulas_tex/formula_339.tex}然后分别删除这两个节点
最后在点集中会剩下\inlineformula{./assets/formulas_tex/formula_340.tex}个点给他们连边
虽然转化方法并不重要 但我们仍给出转化的代码
它有如下的性质
1 \inlineformula{./assets/formulas_tex/formula_341.tex}序列和无根树一一对应
2 度数为\inlineformula{./assets/formulas_tex/formula_342.tex}的点会在prufer序列中出现\inlineformula{./assets/formulas_tex/formula_343.tex}次
3 \inlineformula{./assets/formulas_tex/formula_344.tex}个点的树的\inlineformula{./assets/formulas_tex/formula_345.tex}序列长度为\inlineformula{./assets/formulas_tex/formula_346.tex}
当我们成功把无根树映射成一个序列时 我们将会得到很多好的计数方法
1 对于给定度数序列\inlineformula{./assets/formulas_tex/formula_347.tex}的无根树它的情况数显然为 (可重元素排列计数)
\inlineformula{./assets/formulas_tex/formula_348.tex}
2 对于给定部分度数序列的无根树它的情况数
我们记有度数限制的点有\inlineformula{./assets/formulas_tex/formula_349.tex}个 ，\inlineformula{./assets/formulas_tex/formula_350.tex}=\inlineformula{./assets/formulas_tex/formula_351.tex}
那么方案数为\inlineformula{./assets/formulas_tex/formula_352.tex}}
            \code{src/number_theory/树计数_6.cpp}
        \subsubsection{Cayley定理}
            \txt{1 对于\inlineformula{./assets/formulas_tex/formula_353.tex}个节点的完全图的生成树个数为\inlineformula{./assets/formulas_tex/formula_354.tex}(结论是显然的生成树的prufer序列长为\inlineformula{./assets/formulas_tex/formula_355.tex} 同时\inlineformula{./assets/formulas_tex/formula_356.tex}序列的每个点都有\inlineformula{./assets/formulas_tex/formula_357.tex}种可能性)
2 给定 \inlineformula{./assets/formulas_tex/formula_358.tex}个点（有标号）其生成的无根树的个数为\inlineformula{./assets/formulas_tex/formula_359.tex}(实际上该条件和给定完全图类似)
3给定\inlineformula{./assets/formulas_tex/formula_360.tex}个点（有标号）其生成的有根树的个数为\inlineformula{./assets/formulas_tex/formula_361.tex}(结论是显然的对于\inlineformula{./assets/formulas_tex/formula_362.tex}中的每种情况我们的根节点都有\displayformula{./assets/formulas_tex/formula_363.tex}种选择)}
        \subsubsection{**扩展Cayley定理**}
            \txt{由\inlineformula{./assets/formulas_tex/formula_364.tex}序列推出，给定一个有\inlineformula{./assets/formulas_tex/formula_365.tex}个点\inlineformula{./assets/formulas_tex/formula_366.tex}个联通块的无向图，第\inlineformula{./assets/formulas_tex/formula_367.tex}个联通块的点数为\inlineformula{./assets/formulas_tex/formula_368.tex}，那么要加\inlineformula{./assets/formulas_tex/formula_369.tex}条边使该图联通，方案数为\inlineformula{./assets/formulas_tex/formula_370.tex}}
        \subsubsection{**矩阵树定理**}
            \txt{即给定一个无向图（或有向图）求它的生成树个数（可能是外向树或内向树）
\inlineformula{./assets/formulas_tex/formula_371.tex}求无向图的生成树个数
\inlineformula{./assets/formulas_tex/formula_372.tex}给出有向图和其中的一个点求以这个点为根的生成外向树个数  外向则是由根指向树
\inlineformula{./assets/formulas_tex/formula_373.tex}给出有向图和其中的一个点求以这个点为根的生成内向树的个数 内向则是由树指向根
第一种
所求矩阵\inlineformula{./assets/formulas_tex/formula_374.tex}对角线上为每个点的度数为正数
非对角线的元素\inlineformula{./assets/formulas_tex/formula_375.tex}为\inlineformula{./assets/formulas_tex/formula_376.tex}与\inlineformula{./assets/formulas_tex/formula_377.tex}之间边的数量为负数
去掉最后一行和最后一列求行列式的值
第二种
所求矩阵\inlineformula{./assets/formulas_tex/formula_378.tex}对角线上为每个点的入度为正数
非对角线的元素\inlineformula{./assets/formulas_tex/formula_379.tex}为\inlineformula{./assets/formulas_tex/formula_380.tex}指向\inlineformula{./assets/formulas_tex/formula_381.tex}的边的数量为负数
去掉根所在的行和列求行列式的值
第三种
所求矩阵\inlineformula{./assets/formulas_tex/formula_382.tex}对角线上为每个点的出度为正数
非对角线的元素\inlineformula{./assets/formulas_tex/formula_383.tex}为\displayformula{./assets/formulas_tex/formula_384.tex}指向\inlineformula{./assets/formulas_tex/formula_385.tex}的边的数量为负数
去掉根所在的行和列求行列式的值
对于行列式的求值我们需要配合辗转相减高斯消元
我们需要特别注意的是第一种情况要去掉任意一行一列求值只需要\inlineformula{./assets/formulas_tex/formula_386.tex}自己减一即可
第二种和第三种情况要求删去根所在的行或列也需要对代码进行改动
特别的如果我们要求所有生成树的所有权值和我们只需要把边权\inlineformula{./assets/formulas_tex/formula_387.tex}当做有\inlineformula{./assets/formulas_tex/formula_388.tex}条重边即可体现在矩阵上
根据乘法原理,对于某种生成树的形态,其贡献为每条边重的次数的乘积。
如果把重边次数理解成权值的话,那么矩阵树定理求的就是所有生成树边权乘积的总和。}
            \code{src/number_theory/树计数_2_5.cpp}

    \subsection{Randomized algorithms | 随机化算法}
        \subsubsection{KP算法}
            \txt{该算法是一个能求出高度近似解的贪心算法，能解决的问题是，把一个序列分成两个序列，满足它们和的差值尽量小。
执行流程：不断取出序列中的最大值和次大值，并把他们的差值重新插入序列中。重复该操作直至序列只剩一个元素。那么该元素即为近似解。
如果我们时间允许执行多次，我们不妨每次随机取出一个数字后做该过程，再把取出的数字加进去，这样有概率能随机到最优解。}
        \subsubsection{MinHash}
            \txt{用于比较两个集合的相似度，我们定义两个集合的\inlineformula{./assets/formulas_tex/formula_389.tex}和\inlineformula{./assets/formulas_tex/formula_390.tex}的 Jaccard 相似度记为\inlineformula{./assets/formulas_tex/formula_391.tex},即交集大小除以并集大小，但这样比较的效率很低，我们引入\inlineformula{./assets/formulas_tex/formula_392.tex}技术，我们不妨定义\inlineformula{./assets/formulas_tex/formula_393.tex}个不同的哈希函数，对于每个哈希函数，我们对集合\inlineformula{./assets/formulas_tex/formula_394.tex}中的每个数进行哈希，然后得到其中最小的数字，作为该哈希函数下\inlineformula{./assets/formulas_tex/formula_395.tex}集合的标准值，这样我们可以得到一个长度为\inlineformula{./assets/formulas_tex/formula_396.tex}的向量作为新的\inlineformula{./assets/formulas_tex/formula_397.tex}序列。同理得到\inlineformula{./assets/formulas_tex/formula_398.tex}的新序列。
定理：两个集合的 MinHash 值相同的概率，等于它们的 Jaccard 相似度
这样我们就可以通过比较长度为\inlineformula{./assets/formulas_tex/formula_399.tex}的序列判断两个序列的相似程度。}
        \subsubsection{模拟退火}
            \txt{即随机化找最优解的贪心方法
执行流程是在现有解的基础上进行扰动，但在劣解情况下也有概率进行转移，这样做是为了防止跳入局部最优解中。
步骤一：确定初始温度\displayformula{./assets/formulas_tex/formula_400.tex}和降温系数\displayformula{./assets/formulas_tex/formula_401.tex}以及最低温\displayformula{./assets/formulas_tex/formula_402.tex}，当温度\inlineformula{./assets/formulas_tex/formula_403.tex}时结束算法，并且每一轮过程都让当前温度\inlineformula{./assets/formulas_tex/formula_404.tex}变成\inlineformula{./assets/formulas_tex/formula_405.tex},我们一般取\inlineformula{./assets/formulas_tex/formula_406.tex}为\inlineformula{./assets/formulas_tex/formula_407.tex}。
步骤二：选择一个新解并决定是否接受新解，若答案更优直接接受，或者当一个随机数\inlineformula{./assets/formulas_tex/formula_408.tex},其中\inlineformula{./assets/formulas_tex/formula_409.tex}为两个解的差值
步骤三：选择一个初始点开始执行模拟退火算法即可。
不难发现当\displayformula{./assets/formulas_tex/formula_410.tex}不断趋近于\displayformula{./assets/formulas_tex/formula_411.tex}时，如果新的解不够优秀，接受新解第二个条件变的不可接受，此时会稳定寻找周围的最优解。
初始点的选择：我们尽可能贪心选择一个接近答案位置的解即可。
新解的选择：
我们分不同问题考虑
计算几何连续空间最优化问题：
我们考虑把坐标偏移量和温度挂钩，即使用温度*随机数的方式来选择，这样可以保证在算法初期有足够的探索性，后期又有足够的找精细最优解的能力。
最优顺序/最优选择/最优分组：
方法有很多种，可以随机交换两个位置，也可以翻转某一段区间等等，或者改变某个未知的状态等等。
\# 多项式}
        \subsubsection{拉格朗日插值}
            \txt{根据拉格朗日定理，我们知道,\inlineformula{./assets/formulas_tex/formula_412.tex}对点值\inlineformula{./assets/formulas_tex/formula_413.tex}可以唯一确定一个\inlineformula{./assets/formulas_tex/formula_414.tex}次的多项式，拉格朗日插值即是求解该多项式的过程
用多项式来讲，就是把点值表示法转为系数表示法
其核心思想是考虑点值的独立性然后挨个合并，和\inlineformula{./assets/formulas_tex/formula_415.tex}的思想类似，因此我们想到构造\inlineformula{./assets/formulas_tex/formula_416.tex}个多项式\inlineformula{./assets/formulas_tex/formula_417.tex},满足\inlineformula{./assets/formulas_tex/formula_418.tex},之后最终答案多项式\inlineformula{./assets/formulas_tex/formula_419.tex}
我们有一种显然的构造方法\inlineformula{./assets/formulas_tex/formula_420.tex}然后我们得到了最终的\inlineformula{./assets/formulas_tex/formula_421.tex}
我们发现如果点值是离散数点的话，我们的复杂度将是\inlineformula{./assets/formulas_tex/formula_422.tex}的
但是如果给定点值的横坐标是连续数点的话，我们有单点快速插值的办法
对于分子，我们可以预处理前缀后缀积快速计算，对于分母，我们发现是一个阶乘形式，预处理即可,下面是单点快速插值代码}
            \code{src/number_theory/随机化算法_25.cpp}
        \subsubsection{原根表}
            \code{src/number_theory/随机化算法_2_5.cpp}
        \subsubsection{多项式全家桶}
            \code{src/number_theory/随机化算法_3_5.cpp}
        \subsubsection{集合幂级数全家桶}
            \code{src/number_theory/随机化算法_4_5.cpp}
        \subsubsection{分治NTT}
            \txt{给定一个序列 \inlineformula{./assets/formulas_tex/formula_423.tex} 求序列\displayformula{./assets/formulas_tex/formula_424.tex} 其中 \inlineformula{./assets/formulas_tex/formula_425.tex} 我们将给定边界\inlineformula{./assets/formulas_tex/formula_426.tex}
我们不妨设 \inlineformula{./assets/formulas_tex/formula_427.tex} 设\inlineformula{./assets/formulas_tex/formula_428.tex}
那么\inlineformula{./assets/formulas_tex/formula_429.tex}
因此 \inlineformula{./assets/formulas_tex/formula_430.tex} 需要注意的是把\displayformula{./assets/formulas_tex/formula_431.tex}}
        \subsubsection{Chirp-Z-Transform}
            \txt{假设我们有多项式\inlineformula{./assets/formulas_tex/formula_432.tex}，我们有常数\inlineformula{./assets/formulas_tex/formula_433.tex}和\inlineformula{./assets/formulas_tex/formula_434.tex}，我们希望求出
\inlineformula{./assets/formulas_tex/formula_435.tex}，暴力复杂度为\inlineformula{./assets/formulas_tex/formula_436.tex},我们不妨设\inlineformula{./assets/formulas_tex/formula_437.tex}
代入可知\inlineformula{./assets/formulas_tex/formula_438.tex},我们知道\inlineformula{./assets/formulas_tex/formula_439.tex}
化简以后可以定义\inlineformula{./assets/formulas_tex/formula_440.tex}，\inlineformula{./assets/formulas_tex/formula_441.tex}
注意\inlineformula{./assets/formulas_tex/formula_442.tex}数组大小为\inlineformula{./assets/formulas_tex/formula_443.tex}
我们有\inlineformula{./assets/formulas_tex/formula_444.tex}
也就是\inlineformula{./assets/formulas_tex/formula_445.tex}}
            \code{src/number_theory/随机化算法_5_5.cpp}
        \subsubsection{高精度运算}
            \txt{实现了高精度加法，乘法和减法。
其中高精度乘法分\inlineformula{./assets/formulas_tex/formula_446.tex}的\inlineformula{./assets/formulas_tex/formula_447.tex}优化版本和\inlineformula{./assets/formulas_tex/formula_448.tex}的暴力版本}
            \code{src/number_theory/随机化算法_6_5.cpp}
        \subsubsection{乘法卷积}
            \txt{我们考虑把乘法转化为普通的加法 我们不难想到在指数运算中乘法可以转为加法
但是如果\inlineformula{./assets/formulas_tex/formula_449.tex}，那么\inlineformula{./assets/formulas_tex/formula_450.tex} 那么\inlineformula{./assets/formulas_tex/formula_451.tex}
于是我们需要一个\inlineformula{./assets/formulas_tex/formula_452.tex}意义下的指数和对数映射 且要求是在整数域上的一 一映射
我们不难想到原根 对于质数\inlineformula{./assets/formulas_tex/formula_453.tex}的原根\inlineformula{./assets/formulas_tex/formula_454.tex},有\inlineformula{./assets/formulas_tex/formula_455.tex}在\inlineformula{./assets/formulas_tex/formula_456.tex}意义下两两不同
于是我们把原来的数转化为对数 即\inlineformula{./assets/formulas_tex/formula_457.tex}我们转化为对数的加法即可
请注意我们在执行卷积的时候要循环卷积，特别的，在本算法中对于模数为\inlineformula{./assets/formulas_tex/formula_458.tex}应该是对\inlineformula{./assets/formulas_tex/formula_459.tex}进行循环卷积，因为转化之后模\inlineformula{./assets/formulas_tex/formula_460.tex}乘法群只有\inlineformula{./assets/formulas_tex/formula_461.tex}个元素。
如果我们需要原本乘积为\inlineformula{./assets/formulas_tex/formula_462.tex}的位置的值，那我们应该访问转化后的\inlineformula{./assets/formulas_tex/formula_463.tex}位置，其中\inlineformula{./assets/formulas_tex/formula_464.tex}是我们建立的映射关系。}
            \code{src/number_theory/随机化算法_7_5.cpp}
        \subsubsection{循环卷积}
            \txt{也就是说我们需要的值在\inlineformula{./assets/formulas_tex/formula_465.tex}意义下 我们需要把 \inlineformula{./assets/formulas_tex/formula_466.tex}位置的值加到\inlineformula{./assets/formulas_tex/formula_467.tex}的对应位置}
        \subsubsection{NTT优化字符串通配符匹配}
            \txt{我们先讲常规的\inlineformula{./assets/formulas_tex/formula_468.tex}怎么用\inlineformula{./assets/formulas_tex/formula_469.tex}匹配
我们设长度为\inlineformula{./assets/formulas_tex/formula_470.tex}的模式串\inlineformula{./assets/formulas_tex/formula_471.tex},长度为\inlineformula{./assets/formulas_tex/formula_472.tex}的文本串\inlineformula{./assets/formulas_tex/formula_473.tex},定义匹配函数\inlineformula{./assets/formulas_tex/formula_474.tex}
再定义完美匹配函数\inlineformula{./assets/formulas_tex/formula_475.tex}，根据定义不难发现没如果某个\inlineformula{./assets/formulas_tex/formula_476.tex}满足\inlineformula{./assets/formulas_tex/formula_477.tex}那么我们称\inlineformula{./assets/formulas_tex/formula_478.tex}串以\inlineformula{./assets/formulas_tex/formula_479.tex}结尾的长度为\inlineformula{./assets/formulas_tex/formula_480.tex}的串和\inlineformula{./assets/formulas_tex/formula_481.tex}匹配
我们把式子展开得到\inlineformula{./assets/formulas_tex/formula_482.tex}
再把a串翻转得到
\inlineformula{./assets/formulas_tex/formula_483.tex}
再用\inlineformula{./assets/formulas_tex/formula_484.tex}代替\inlineformula{./assets/formulas_tex/formula_485.tex}
\inlineformula{./assets/formulas_tex/formula_486.tex}
我们展开
\inlineformula{./assets/formulas_tex/formula_487.tex}
我们发现第一项为定值，第二项可以前缀和处理，
第三项为卷积 于是我们在\inlineformula{./assets/formulas_tex/formula_488.tex}时间内解决了这个问题
然后我们引进通配符\inlineformula{./assets/formulas_tex/formula_489.tex}事实上我们只需要让通配符的值为\inlineformula{./assets/formulas_tex/formula_490.tex}且改变匹配函数的定义即可
\inlineformula{./assets/formulas_tex/formula_491.tex}
我们再次展开
\displayformula{./assets/formulas_tex/formula_492.tex}
变成了三个卷积，但无伤大雅}
            \code{src/number_theory/随机化算法_8_5.cpp}
        \subsubsection{NTT优化有限字符集的匹配}
            \txt{当我们的字符集数量有限且时间充裕时，我们可以考虑每次只匹配一个字符，然后对每个字符都跑一次匹配，累加起来便是答案。
这样做的好处是我们可以得到每个位置的匹配情况，即能有多少个位置匹配上.
具体的说，假设我们匹配字符\displayformula{./assets/formulas_tex/formula_493.tex},我们把两个串对应位置的权值设为\displayformula{./assets/formulas_tex/formula_494.tex}然后直接把权值卷起来即可，注意我们仍然要翻转字符串\inlineformula{./assets/formulas_tex/formula_495.tex}
复杂度\inlineformula{./assets/formulas_tex/formula_496.tex}
\# 博弈论
主要分为公平组合博弈和不公平游戏
分别有不同的考虑方法}
        \subsubsection{公平组合游戏}
            \txt{即双方的可行操作集合相同
对于公平组合游戏，我们可以考虑\displayformula{./assets/formulas_tex/formula_497.tex}定理，把组合游戏拆分成独立的小游戏，算出每个小游戏的\inlineformula{./assets/formulas_tex/formula_498.tex}函数值，它们的异或和为组合游戏的\inlineformula{./assets/formulas_tex/formula_499.tex}值，也称为\inlineformula{./assets/formulas_tex/formula_500.tex}和，其中\inlineformula{./assets/formulas_tex/formula_501.tex}，其中\inlineformula{./assets/formulas_tex/formula_502.tex}表示\inlineformula{./assets/formulas_tex/formula_503.tex}能到达的状态。
如果不容易求\inlineformula{./assets/formulas_tex/formula_504.tex}函数或者不满足\inlineformula{./assets/formulas_tex/formula_505.tex}函数的条件，我们还可以考虑记忆化搜索的方法暴力求解，即后继态有必胜态的为必胜态，全为必败态的为必败态。}
            \code{src/number_theory/随机化算法_9_5.cpp}
        \subsubsection{公平组合游戏下的经典模型}
        \subsubsection{Nim游戏}
            \txt{给定\inlineformula{./assets/formulas_tex/formula_506.tex}堆式子，第\inlineformula{./assets/formulas_tex/formula_507.tex}堆石子有\inlineformula{./assets/formulas_tex/formula_508.tex}个，两名玩家轮流行动，每次任选一堆取走任意多个石子，不能不取，取走最后一颗石子者获胜。
结论：定义\inlineformula{./assets/formulas_tex/formula_509.tex}和为\inlineformula{./assets/formulas_tex/formula_510.tex}，若\inlineformula{./assets/formulas_tex/formula_511.tex}和为\inlineformula{./assets/formulas_tex/formula_512.tex}，那么先手必败，否则先手必胜。
证明：
若当前\inlineformula{./assets/formulas_tex/formula_513.tex}和不为\inlineformula{./assets/formulas_tex/formula_514.tex}，我一定可以使得下一步\inlineformula{./assets/formulas_tex/formula_515.tex}和为\inlineformula{./assets/formulas_tex/formula_516.tex},也即必胜态以后一定有必败态。
若当前\inlineformula{./assets/formulas_tex/formula_517.tex}和为\inlineformula{./assets/formulas_tex/formula_518.tex}，我一定找不到一种方法使得下一步\inlineformula{./assets/formulas_tex/formula_519.tex}和仍然为\inlineformula{./assets/formulas_tex/formula_520.tex}，即必败态后都是必胜态}
        \subsubsection{巴什博奕}
            \txt{有一堆石子数量为\inlineformula{./assets/formulas_tex/formula_521.tex},两名玩家轮流取，拿的个数范围是\displayformula{./assets/formulas_tex/formula_522.tex}，取走最后一颗石子者获胜。
结论：若\inlineformula{./assets/formulas_tex/formula_523.tex}，则先手必败。反之先手必胜。
\inlineformula{./assets/formulas_tex/formula_524.tex}
证明:
若\inlineformula{./assets/formulas_tex/formula_525.tex}显然先手必胜，当\inlineformula{./assets/formulas_tex/formula_526.tex}时显然先手必败，当\inlineformula{./assets/formulas_tex/formula_527.tex}时，如果\inlineformula{./assets/formulas_tex/formula_528.tex}是\inlineformula{./assets/formulas_tex/formula_529.tex}的倍数，无论先手如何操作，都无法把必败态留给对方。反之，后手一定有办法再次使得\inlineformula{./assets/formulas_tex/formula_530.tex}，也即把必败态留给先手。}
        \subsubsection{威佐夫博弈}
            \txt{有两堆石子，两人轮流取石子，每次可以在一堆中取，也可以在两堆中取走相同数量的石子。取走最后一个石子者获胜。
结论：记\inlineformula{./assets/formulas_tex/formula_531.tex}为两堆石子个数，若\inlineformula{./assets/formulas_tex/formula_532.tex}，那么先手必败，否则先手必胜。
证明略。}
        \subsubsection{斐波那契博弈}
            \txt{有一堆数量为\inlineformula{./assets/formulas_tex/formula_533.tex}的石子，两人轮流取，先手不能一次取完。如果前一个人取了\inlineformula{./assets/formulas_tex/formula_534.tex}个石子，那么我取的石子数量范围为\inlineformula{./assets/formulas_tex/formula_535.tex}，取走最后一个石子者获胜。
结论:若\inlineformula{./assets/formulas_tex/formula_536.tex}为斐波那契数则先手必败，否则先手必胜。}
        \subsubsection{广义斐波那契博弈}
            \txt{即如果前一个人取了\inlineformula{./assets/formulas_tex/formula_537.tex},我能取的范围变为\inlineformula{./assets/formulas_tex/formula_538.tex}，此时我们有广义的斐波那契数列\displayformula{./assets/formulas_tex/formula_539.tex}，即\inlineformula{./assets/formulas_tex/formula_540.tex}如果石子数量在该数列中，则先手必败。否则必胜}
        \subsubsection{Nim-K 博弈}
            \txt{在\inlineformula{./assets/formulas_tex/formula_542.tex}的基础上，每次可以从\inlineformula{./assets/formulas_tex/formula_543.tex}堆石子取，每堆取的石子数量可以不同。
一堆石子的\inlineformula{./assets/formulas_tex/formula_544.tex}值为石子数，对每一个二进制位单独考虑，求每一个二进制位为\displayformula{./assets/formulas_tex/formula_545.tex}的堆数，然后\inlineformula{./assets/formulas_tex/formula_546.tex}，如果每个二进制的结果均为\inlineformula{./assets/formulas_tex/formula_547.tex}则必败，否则必胜}
        \subsubsection{Anti-Nim 博弈}
            \txt{即获胜条件改为不能取的一方获胜
必胜：\inlineformula{./assets/formulas_tex/formula_549.tex}且至少有一堆石子数\inlineformula{./assets/formulas_tex/formula_550.tex},\inlineformula{./assets/formulas_tex/formula_551.tex}且每堆石子个数不超过1
其余均必败}
        \subsubsection{Anti-SG 博弈}
            \txt{获胜条件改为不能移动的一方获胜
必胜：\inlineformula{./assets/formulas_tex/formula_553.tex}且至少有一个游戏满足\inlineformula{./assets/formulas_tex/formula_554.tex},\inlineformula{./assets/formulas_tex/formula_555.tex}且每个游戏都满足\displayformula{./assets/formulas_tex/formula_556.tex}
其余均必败}
        \subsubsection{Staircase-Nim 博弈}
            \txt{阶梯博弈，每次可以从一个阶梯上拿掉任意数量石子放到下一个阶梯，不能操作者输
\inlineformula{./assets/formulas_tex/formula_558.tex}函数为奇数阶梯上的石子的异或和，如果移动偶数层的石子到奇数层，对手一定能把这些石子继续移动到偶数层，使其\inlineformula{./assets/formulas_tex/formula_559.tex}函数不变。
\inlineformula{./assets/formulas_tex/formula_560.tex}则必败。否则必胜}
        \subsubsection{Lasker's-Nim 博弈}
            \txt{\displayformula{./assets/formulas_tex/formula_562.tex}堆石子，每次可以从一堆中取任意个石子，或者至少选择某堆至少为\displayformula{./assets/formulas_tex/formula_563.tex}的石子，分成两堆非空石子。
\displayformula{./assets/formulas_tex/formula_564.tex}
对于\displayformula{./assets/formulas_tex/formula_565.tex},满足\displayformula{./assets/formulas_tex/formula_566.tex},\displayformula{./assets/formulas_tex/formula_567.tex}
\displayformula{./assets/formulas_tex/formula_568.tex}}
        \subsubsection{树上删边博弈}
            \txt{给一个\displayformula{./assets/formulas_tex/formula_569.tex}个点的有根树，每次可以删掉一个子树。
叶子结点的\displayformula{./assets/formulas_tex/formula_570.tex}为\displayformula{./assets/formulas_tex/formula_571.tex}，非叶子节点的\displayformula{./assets/formulas_tex/formula_572.tex}为所有子节点的\displayformula{./assets/formulas_tex/formula_573.tex}的异或和}
        \subsubsection{无向图删边游戏}
            \txt{注意此时也需要一个参考点（也就是根），删去一条边后，不属于参考点所在的联通块会被舍去。
对于自环的处理
（把它连向一个新的点即可，或者让它自己的函数值异或\inlineformula{./assets/formulas_tex/formula_574.tex}）
我们考虑\displayformula{./assets/formulas_tex/formula_575.tex}缩点，每个\displayformula{./assets/formulas_tex/formula_576.tex}的\displayformula{./assets/formulas_tex/formula_577.tex}值取决于\displayformula{./assets/formulas_tex/formula_578.tex}内边的数量，如有偶数条边则为\displayformula{./assets/formulas_tex/formula_579.tex},否则为\displayformula{./assets/formulas_tex/formula_580.tex}
然后套用树上删边博弈即可。}
        \subsubsection{一般图上博弈}
            \txt{关于二分图上博弈的内容在二分图中有介绍
给定无向图, 先手从起点开始走, 只能走相邻且未走过的点, 无法移动者输.
对该图求最大匹配, 若起点不一定在最大匹配中则先手必败, 否则先手必胜.}
        \subsubsection{翻硬币游戏}
            \txt{\displayformula{./assets/formulas_tex/formula_581.tex}枚硬币排成一排，有的正面朝上，有的反面朝上，给定一些翻硬币的约束（比如只能翻\displayformula{./assets/formulas_tex/formula_582.tex}枚，比如只能翻连续几枚），但所翻的硬币中，最右边的必须是从正面翻到反面，不能行动者输。
有如下结论：局面的\displayformula{./assets/formulas_tex/formula_583.tex}值等于局面中每个正面朝上的棋子单一存在时的\displayformula{./assets/formulas_tex/formula_584.tex}值的异或和。（即只有该棋子正面，其他都是反面，根据它的位置\inlineformula{./assets/formulas_tex/formula_585.tex}有一个\inlineformula{./assets/formulas_tex/formula_586.tex}）
情况\displayformula{./assets/formulas_tex/formula_587.tex}，每次只能翻一个
\displayformula{./assets/formulas_tex/formula_588.tex}
情况\inlineformula{./assets/formulas_tex/formula_589.tex}，每次可以翻一个或两个(不要求相邻)
\inlineformula{./assets/formulas_tex/formula_590.tex}
情况\inlineformula{./assets/formulas_tex/formula_591.tex}，每次必须翻两个，且两个的距离在\inlineformula{./assets/formulas_tex/formula_592.tex}中，等价于巴什博奕
\inlineformula{./assets/formulas_tex/formula_593.tex}
情况\inlineformula{./assets/formulas_tex/formula_594.tex}：每次必须翻连续的\inlineformula{./assets/formulas_tex/formula_595.tex}个硬币
\displayformula{./assets/formulas_tex/formula_596.tex},其余为\displayformula{./assets/formulas_tex/formula_597.tex}
情况\displayformula{./assets/formulas_tex/formula_598.tex}每次可以翻任意长度的连续一段硬币，\displayformula{./assets/formulas_tex/formula_599.tex}}
        \subsubsection{不公平游戏}
            \txt{两名玩家轮流操作的博弈，每个局面下 双方可行的操作集合可能不同，或者某些动作只允许某一方使用。
换句话说，玩家的可行动作不对称，我们称为不公平游戏。}
        \subsubsection{超现实数 (Surreal Numbers)}
            \txt{超现实数类似于\displayformula{./assets/formulas_tex/formula_600.tex}函数，用于解决不公平游戏问题
超现实数的不同游戏间的组合方法是相加
我们不妨用\displayformula{./assets/formulas_tex/formula_601.tex}来表示
如果\displayformula{./assets/formulas_tex/formula_602.tex} 那么左玩家必胜
如果\displayformula{./assets/formulas_tex/formula_603.tex}那么右玩家必胜
如果\displayformula{./assets/formulas_tex/formula_604.tex}那么后手必胜
单个游戏的值\displayformula{./assets/formulas_tex/formula_605.tex}被定义为\displayformula{./assets/formulas_tex/formula_606.tex},其中\inlineformula{./assets/formulas_tex/formula_607.tex}是左玩家所有可能移动后到达的局面集合，\inlineformula{./assets/formulas_tex/formula_608.tex}同理。}
        \subsubsection{Blue-Red Hackenbush 博弈}
            \txt{有若干个字符串仅有\inlineformula{./assets/formulas_tex/formula_610.tex}和\inlineformula{./assets/formulas_tex/formula_611.tex}组成，小\inlineformula{./assets/formulas_tex/formula_612.tex} 回合只能拿W，小\inlineformula{./assets/formulas_tex/formula_613.tex}回合只能拿\inlineformula{./assets/formulas_tex/formula_614.tex}，每次拿走一个字符后其后缀也会消失，最先不能操作者输。
首先这是一个不公平游戏，因为操作域不同，因此我们考虑对每个串计算它的\inlineformula{./assets/formulas_tex/formula_615.tex}并求和，再根据上面结论判断即可。
在本题中可以证明，一个串的\inlineformula{./assets/formulas_tex/formula_616.tex}值为\inlineformula{./assets/formulas_tex/formula_617.tex}，
其中\inlineformula{./assets/formulas_tex/formula_618.tex}。}
        \subsubsection{Termites}
            \txt{一个序列，两个人轮流从最左或者最右拿一个数，谁总和大谁胜。
转化成求最大差值，对于连续的三个数\inlineformula{./assets/formulas_tex/formula_620.tex}，如果\inlineformula{./assets/formulas_tex/formula_621.tex} 且\inlineformula{./assets/formulas_tex/formula_622.tex}，那么可以将这三个数合并为\inlineformula{./assets/formulas_tex/formula_623.tex}。从左往右维护一个栈，每次新加入数的时候不断合并栈顶三个元素，那么最后栈是先递减再递增，此时贪心拿两侧最大的数即为最佳策略。}
        \subsubsection{Game of Sorting}
            \txt{给出一列数，两人轮流从最左或者最右拿走一个数，如果目前轮到\inlineformula{./assets/formulas_tex/formula_625.tex}操作，这个序列单调不下降或者单调不上升，那么\inlineformula{./assets/formulas_tex/formula_626.tex}输。
每次询问在某个区间\inlineformula{./assets/formulas_tex/formula_627.tex} 玩游戏的结果。
我们考虑什么样的情况能判断结果
\inlineformula{./assets/formulas_tex/formula_628.tex}对于一个区间\inlineformula{./assets/formulas_tex/formula_629.tex}，若它是单调的，则显然先手必败。
\inlineformula{./assets/formulas_tex/formula_630.tex}若去掉\inlineformula{./assets/formulas_tex/formula_631.tex} 或者\inlineformula{./assets/formulas_tex/formula_632.tex} 后变成了单调的，则显然先手必胜。
\inlineformula{./assets/formulas_tex/formula_633.tex}若\inlineformula{./assets/formulas_tex/formula_634.tex} 是单调的，那么同理先手必败。
\inlineformula{./assets/formulas_tex/formula_635.tex}若\inlineformula{./assets/formulas_tex/formula_636.tex} 是单调的，那么先手若取\inlineformula{./assets/formulas_tex/formula_637.tex}会导致后手必胜，后手同理，故两人会一直取\inlineformula{./assets/formulas_tex/formula_638.tex}直到出现上面第一种或者第二种情况，可以根据奇偶性判断。
\inlineformula{./assets/formulas_tex/formula_639.tex}同理可以得出\inlineformula{./assets/formulas_tex/formula_640.tex} 是单调的情形的处理方法。
除此之外，有\inlineformula{./assets/formulas_tex/formula_641.tex}的答案等于\inlineformula{./assets/formulas_tex/formula_642.tex} 的答案，二分找到最小的k 满足\inlineformula{./assets/formulas_tex/formula_643.tex}可以由上面方法直接判断出胜负即可。}
        \subsubsection{Xormites}
            \txt{给出一列数，两人轮流从最左或者最右拿走一个数，最后得分是拿过的数的异或和，谁高
谁赢，预测结果
\inlineformula{./assets/formulas_tex/formula_645.tex}首先如果所有数字异或和为\inlineformula{./assets/formulas_tex/formula_646.tex}，一定平局
\inlineformula{./assets/formulas_tex/formula_647.tex}在除去\inlineformula{./assets/formulas_tex/formula_648.tex}的情况下，一定不会平局
\inlineformula{./assets/formulas_tex/formula_649.tex}当\inlineformula{./assets/formulas_tex/formula_650.tex}为偶数时，先手必胜，我们不妨给\inlineformula{./assets/formulas_tex/formula_651.tex}个格子相邻黑白染色，不难发现先手可以选择拿所有的白色格子或者黑色格子，这两个选大的拿即可。
\inlineformula{./assets/formulas_tex/formula_652.tex}我们不妨考虑异或和\inlineformula{./assets/formulas_tex/formula_653.tex}的最高二进制位\inlineformula{./assets/formulas_tex/formula_654.tex},那么我们根据这一个数字是否有\inlineformula{./assets/formulas_tex/formula_655.tex}这一位分为\inlineformula{./assets/formulas_tex/formula_656.tex}或\inlineformula{./assets/formulas_tex/formula_657.tex}，那么拿到奇数个\inlineformula{./assets/formulas_tex/formula_658.tex}则必胜。如果\inlineformula{./assets/formulas_tex/formula_659.tex}，那么后手模仿方案\inlineformula{./assets/formulas_tex/formula_660.tex}即可取胜。否则先手一定会拿\inlineformula{./assets/formulas_tex/formula_661.tex},先手拿\inlineformula{./assets/formulas_tex/formula_662.tex}以后，必须模仿后手行动，即后手拿什么则先手拿什么，假如先手后手拿的不同，那么此时先手和后手手里的数字相同且轮到后手，后手仍然可以通过\inlineformula{./assets/formulas_tex/formula_663.tex}必胜，因此我们判断能否模仿即可
\inlineformula{./assets/formulas_tex/formula_664.tex}如何判断能否模仿\inlineformula{./assets/formulas_tex/formula_665.tex}，我们先不断删去首尾相同的数字，然后判断接下来两两是否相同即可
\# 线性代数}
        \subsubsection{矩阵与矩阵快速幂}
            \txt{十分需要注意的是 矩阵的幂由于不满足费马小定理 因此不能取模}
            \code{src/number_theory/随机化算法_11_5.cpp}
        \subsubsection{行列式}
            \txt{类似于高斯消元消成上三角形式，这里我们采用了辗转相减的方法，可以应对\inlineformula{./assets/formulas_tex/formula_666.tex}数为非质数的情况}
            \code{src/number_theory/随机化算法_12_5.cpp}
        \subsubsection{加法高斯消元}
            \txt{我们通过\inlineformula{./assets/formulas_tex/formula_667.tex}的三重循环将增广矩阵消成上三角形式，再从末尾式子开始回代得到答案，其中若系数全都为\inlineformula{./assets/formulas_tex/formula_668.tex}但解不为\inlineformula{./assets/formulas_tex/formula_669.tex}，则说明无解。否则有解，如果存在自由元，那么有无穷多解，否则有唯一解}
            \code{src/number_theory/随机化算法_13_5.cpp}
        \subsubsection{异或高斯消元}
            \txt{复杂度为\inlineformula{./assets/formulas_tex/formula_670.tex}，做法和加法类似，使用\inlineformula{./assets/formulas_tex/formula_671.tex}优化。}
            \code{src/number_theory/随机化算法_14_5.cpp}
        \subsubsection{动态维护异或方程组}
            \txt{维护一个与方程组变量数相同的线性基即可。每插入一条新的异或等式时，
先将其化为 $\sum b_i x_i = v$（按位异或），随后使用线性基消去出现的最高位；
若无法消去则把该向量加入基底。线性基记录每一位的代表向量，就能在对数复杂度内判断方程组的可解性，并支持在线更新。}
        \subsubsection{二维平面上的距离}
            \txt{设点 $A(x_1,y_1)$ 与 $B(x_2,y_2)$，常用距离如下：}
            \txt{$$\operatorname{Manhattan}(A,B) = |x_1-x_2| + |y_1-y_2|$$}
            \txt{$$\operatorname{Chebyshev}(A,B) = \max\big(|x_1-x_2|,\ |y_1-y_2|\big)$$}
            \txt{$$\operatorname{Euclidean}(A,B) = \sqrt{(x_1-x_2)^2 + (y_1-y_2)^2}$$}
            \txt{对坐标作变换 $u = x + y,\ v = x - y$ 可以在曼哈顿坐标系与切比雪夫坐标系之间互相转换，从而简化最近点对或旋转卡壳类问题。}
        \subsubsection{自适应辛普森法}
            \txt{自适应辛普森法的用途主要是高精度数值积分，即计算积分值。}
            \code{src/number_theory/随机化算法_16_5.cpp}
        \subsubsection{平面最近点对}
            \code{src/number_theory/随机化算法_17_5.cpp}
        \subsubsection{最小圆覆盖和最小球覆盖}
            \code{src/number_theory/随机化算法_18_5.cpp}
        \subsubsection{单峰函数与三分}
            \txt{常见单峰函数：
二次函数：其最值在对称轴取得
绝对值函数：
\inlineformula{./assets/formulas_tex/formula_672.tex}是单峰函数，其最值在所有\inlineformula{./assets/formulas_tex/formula_673.tex}的中位数取得
整数三分
实数三分
三分套三分
使用场景往往是面向二元函数\inlineformula{./assets/formulas_tex/formula_674.tex},满足在固定\inlineformula{./assets/formulas_tex/formula_675.tex}时\inlineformula{./assets/formulas_tex/formula_676.tex}随\inlineformula{./assets/formulas_tex/formula_677.tex}变化是单峰的，那么我们对固定的\inlineformula{./assets/formulas_tex/formula_678.tex}，可三分出最优的\inlineformula{./assets/formulas_tex/formula_679.tex}。同时我们定义\inlineformula{./assets/formulas_tex/formula_680.tex}，如果\inlineformula{./assets/formulas_tex/formula_681.tex}也是单峰的，即对于每个\inlineformula{./assets/formulas_tex/formula_682.tex}在取到最优的\inlineformula{./assets/formulas_tex/formula_683.tex}以后仍然保持单峰，那么就可以采用三分套三分。
结论：多个下凸函数的\inlineformula{./assets/formulas_tex/formula_684.tex}仍然是下凸函数，多个上凸函数的\inlineformula{./assets/formulas_tex/formula_685.tex}仍然是上凸函数}
            \code{src/number_theory/随机化算法_20_5.cpp}

\section{Computational geometry | 计算几何.cpp}
    \subsection{几何基础}
        \txt{包含点、线、线段、多边形、凸多边形与圆的常用数据结构及运算。}
        \code{src/computational_geometry/geometry_base.cpp}
    \subsection{圆与多边形面积交}
        \code{src/computational_geometry/circle_polygon_area.cpp}
    \subsection{凸包}
        \code{src/computational_geometry/convex_hull.cpp}
    \subsection{半平面交}
        \code{src/computational_geometry/halfplane_intersection.cpp}
    \subsection{点集形成的最小最大三角形}
        \code{src/computational_geometry/minmax_triangle.cpp}
    \subsection{平面最近点对}
        \code{src/computational_geometry/closest_pair.cpp}
    \subsection{线段扫线判交}
        \code{src/computational_geometry/segment_sweep_intersection.cpp}
    \subsection{多边形面积并}
        \code{src/computational_geometry/polygon_union.cpp}
    \subsection{圆面积并}
        \code{src/computational_geometry/circle_union.cpp}

\section{Misc | 杂项}

    \subsection{Bitset | 暴力滴神}
        \txt{set(x): 将 x 位设置为 1，如果不填则全部为 1}
        \txt{reset(x): 将 x 位设置为 0，如果不填则全部为 0}
        \txt{flip(x): 将 x 位取反，如果不填则全部取反}
        \txt{count(): 返回 1 的个数}
        \txt{size(): 返回位数}
        \txt{any(): 返回是否存在 1}
        \txt{\_Find\_fisrt(): 返回第一个 1 的位置}
        \txt{\_Find\_next(x): 返回 x 之后第一个 1 的位置}

\section{Temp | 临时式子整理}

    \subsection{交换律}
    $p(k)$ 是指标集 $k$ 的一个排列。
    \begin{align*}
        \sum_{i \in k} a_i &= \sum_{i \in p(k)} a_i \\
        a_1 + a_2 + a_3 + a_4 &= a_2 + a_1 + a_3 + a_4
    \end{align*}

    \subsection{结合律}
    \begin{align*}
        \sum_{i=1}^{n} a_i + \sum_{i=1}^{n} b_i = \sum_{i=1}^{n} (a_i + b_i)
    \end{align*}

    \subsection{分配律}
    \begin{align*}
        \sum_{i} c \cdot a_i = c \cdot \sum_{i} a_i
    \end{align*}


\section{Summation Transformation | 和式的变换技巧}
    \subsection{替换条件式}
    把写在求和符号下标位置的条件改为艾弗森括号的形式，使总的求和式化为枚举所有变量的求和式，便于处理。
    \begin{align*}
        \sum_{d|\gcd(i,j)} f(d) = \sum_{d=1}^{\min(i,j)} [d|i][d|j]f(d)
    \end{align*}

    \subsection{替换指标变量}
    \begin{align*}
        &\sum_{i=1}^{n} \sum_{j=1}^{m} [\gcd(i,j) = k] \\
        &= \sum_{i=1}^{\lfloor n/k \rfloor} \sum_{j=1}^{\lfloor m/k \rfloor} [\gcd(ik,jk)=k] \\
        &= \sum_{i=1}^{\lfloor n/k \rfloor} \sum_{j=1}^{\lfloor m/k \rfloor} [\gcd(i,j)=1]
    \end{align*}
    \begin{align*}
        \sum_{i=0}^{n} (-1)^{n-i} f(n-i) = \sum_{i=0}^{n} (-1)^{i} f(i)
    \end{align*}

    \subsection{改变求和次序}
    \subsubsection{基础}
    \begin{align*}
        \sum_i \sum_j f(i,j) 
        &= \sum_i \sum_j f(i,j) [1 \le i \le n] [1 \le j \le i] \\
        &= \sum_j \sum_i f(i,j) [1 \le j \le i] [1 \le i \le n] \\
        &= \sum_j \sum_i f(i,j) [j \le i \le n] [1 \le j \le n] \\
        &= \sum_{j=1}^{n} \sum_{i=j}^{n} f(i,j)
    \end{align*}
    \subsubsection{复杂的改变求和次序}
    为了改变求和次序，有时候我们可能需要通过解析不等式的方式。
    \begin{align*}
        &\sum_{i=1}^{n} \sum_{j=i+1}^{n} f(i,j) \\
        &= \sum_j \sum_i f(i,j) [1 \le i \le j-1] [2 \le j \le n] \\
        &= \sum_{j=2}^{n} \sum_{i=1}^{j-1} f(i,j)
    \end{align*}

    \subsection{分离变量}
    \begin{align*}
        \sum_i \sum_j f(i)g(j) = \left( \sum_i f(i) \right) \left( \sum_j g(j) \right)
    \end{align*}

    \subsection{改变枚举范围}
    \subsubsection{扩大枚举范围}
    扩大范围可以将变量有关的枚举上下界改变为常量，便于进行下一步处理。
    \begin{align*}
        &\sum_{i=1}^{n} \sum_{j=1}^{m} \sum_{d=1}^{\min(i,j)} [d|i][d|j]f(d) \\
        &= \sum_{d=1}^{n} \sum_{i=1}^{n} \sum_{j=1}^{m} [d|i][d|j]f(d)
    \end{align*}

    \subsubsection{缩小枚举范围}
    缩小枚举范围使枚举条件简化，也是一个常见技巧。
    \begin{align*}
        &\sum_{i=1}^{n} \sum_{j=1}^{m} [\gcd(i,j) = k] \\
        &= \sum_{i=1}^{\lfloor n/k \rfloor} \sum_{j=1}^{\lfloor m/k \rfloor} [\gcd(i,j)=1]
    \end{align*}

    \subsection{示例：组合数}
    设组合数 $\binom{n}{k}$ 表示从 $n$ 个不同物品中取出 $k$ 个，不考虑取出的顺序，有多少种方案。
    \begin{align*}
        &\sum_{i=0}^{n-1} \binom{n-1}{i} n^{i+1} + \sum_{i=0}^{n} \binom{n-1}{i} n^i \\
        &= \sum_{i=0}^{n} \binom{n-1}{i} n^{i+1} + \sum_{i=0}^{n} \binom{n-1}{i} n^i \\
        &= \sum_{i=0}^{n} \binom{n-1}{i-1} n^{i} + \sum_{i=0}^{n} \binom{n-1}{i} n^i \\
        &= \sum_{i=0}^{n} \left(\binom{n-1}{i-1} + \binom{n-1}{i}\right) n^i \\
        &= \sum_{i=0}^{n} \binom{n}{i} n^i
    \end{align*}

    \subsection{和式的变换公式}
    \subsubsection{公式 1}
    这个式子似乎没什么用。
    \begin{align*}
        \left(\sum_{i=1}^{n} a_i\right)^2 = \sum_{i=1}^{n} a_i^2 + 2 \sum_{1 \le i < j \le n} a_i a_j
    \end{align*}

    \subsubsection{公式 2}
    这个公式也是需要背的，因为其实用化不了时间复杂度，直接背比较好。证明也很容易：
    \begin{align*}
        &\sum_{1 \le i < j \le n} (a_i - a_j)^2 \\
        &= \sum_{1 \le i < j \le n} (a_i^2 + a_j^2 - 2a_i a_j) \\
        &= \sum_{1 \le i < j \le n} a_i^2 + \sum_{1 \le i < j \le n} a_j^2 - 2\sum_{1 \le i < j \le n} a_i a_j \\
        &= (n-1)\sum_{i=1}^{n} a_i^2 - \left(\left(\sum_{i=1}^{n} a_i\right)^2 - \sum_{i=1}^{n} a_i^2\right) \\
        &= n \sum_{i=1}^{n} a_i^2 - \left(\sum_{i=1}^{n} a_i\right)^2
    \end{align*}

    \subsection{阿贝尔部分求和公式}
    这个式子也是中看不中用。
    \begin{align*}
        &\sum_{i=1}^{n} a_i b_i = \\
        &b_n \sum_{i=1}^{n} a_i - \sum_{i=1}^{n-1} \left(\sum_{j=1}^{i} a_j\right) (b_{i+1} - b_i)
    \end{align*}

\end{multicols}
\end{document}